\documentclass[11pt,a4paper]{article}

%---------- packages ----------
\usepackage[utf8]{inputenc}
\usepackage[T1]{fontenc}
\usepackage{amsmath, amssymb, amsthm, mathtools}
\usepackage{microtype}
\usepackage[margin=1in]{geometry}
\usepackage[colorlinks=true, allcolors=blue]{hyperref}
\usepackage{cleveref}

%---------- theorem environments ----------
\theoremstyle{plain}
\newtheorem{theorem}{Theorem}[section]
\newtheorem{proposition}[theorem]{Proposition}
\newtheorem{lemma}[theorem]{Lemma}
\newtheorem{corollary}[theorem]{Corollary}

\theoremstyle{definition}
\newtheorem{definition}[theorem]{Definition}

\theoremstyle{remark}
\newtheorem{remark}[theorem]{Remark}
\newtheorem*{remark*}{Remark}

%---------- title ----------
\title{A Combinatorial Sieve Lower Bound for Primes in the Family
$m \cdot 2^{k} \cdot 3^{\ell} + 1$\\
{\small preprint, expert review pending}}

\author{Jared Wilder\thanks{\texttt{jared@jaredwilder.com}.
Independent researcher. Drafting process disclosed in
\cref{sec:limitations}.}}

\date{2 June 2026}

\begin{document}

\maketitle

\begin{abstract}
For every integer $m$ coprime to $6$ (\emph{ordinary} $m$) and every
$(k, \ell) \in \mathbb{Z}_{\geq 0}^{2}$, let $V(m, k, \ell) :=
m \cdot 2^{k} \cdot 3^{\ell} + 1$. Let $\pi_{V}(m, D)$ count the
$(k, \ell)$ with $k + \ell \leq D$ such that $V(m, k, \ell)$ is prime.
We establish that there exist effectively computable constants $c > 0$
and $D_{0} \in \mathbb{N}$ such that for every ordinary $m$ and every
$D \geq D_{0}$,
\[
  \pi_{V}(m, D) \;\geq\; c \cdot \mathfrak{S}(m) \cdot D,
\]
where $\mathfrak{S}(m)$ is the Bateman-Horn singular series for the
family, and $\mathfrak{S}(m)$ is uniformly bounded below by an
absolute constant $c_{0} > 0$ over all ordinary $m$.

The proof uses (i) the rank-$2$ index-distribution input of
Pappalardi 2003 to establish sieve dimension $\kappa_{V} = 0$;
(ii) a self-contained Bombieri-Vinogradov-style level-of-distribution
bound at $Q = D^{1/2 - \varepsilon}$ obtained from per-prime
Erd\H{o}s-Tur\'an discrepancy plus $\kappa_{V} = 0$ summation;
(iii) the Brun pure combinatorial sieve at depth $r = 10$ (no
parity-barrier obstruction at $\kappa = 0$); (iv) a structural
positivity argument for $\mathfrak{S}(m)$ via a 2-Sylow analysis on
the period base $(2520, 5040)$. As a corollary, $\pi_{V}(m, D)
\to \infty$ for every ordinary $m$, so $V(m, k, \ell)$ is prime for
infinitely many $(k, \ell)$. This affirms Erd\H{o}s-Graham
Problem~203 (1980).

Our effective $D_{0} \approx 10^{87}$ is far above the empirical
record ($D \leq 22$ for $m \leq 2 \cdot 10^{9}$); closing the gap is
left open.

\textbf{Status:} preprint quality; expert review solicited. Literature
citations are verified at the level of paper existence and theorem
location only; verbatim text checking is incomplete (see
\cref{sec:limitations}).
\end{abstract}

\tableofcontents

\section*{Executive summary}

The Erd\H{o}s-Graham conjecture~\#203 asks whether for every ordinary
$m$ (i.e.\ $\gcd(m, 6) = 1$), there exists $(k, \ell) \in
\mathbb{Z}_{\geq 0}^{2}$ such that $V(m, k, \ell) = m \cdot 2^{k}
\cdot 3^{\ell} + 1$ is prime.

\textbf{Main result.} \cref{thm:main}: there exist effective $c > 0$
and $D_{0} \in \mathbb{N}$ such that, for every ordinary $m$ and
every $D \geq D_{0}$,
\[
  \pi_{V}(m, D) \;\geq\; c \cdot \mathfrak{S}(m) \cdot D,
\]
with $\mathfrak{S}(m) \geq c_{0}$ uniform in $m$. Corollary: $\pi_{V}(m,
D) \to \infty$ as $D \to \infty$, so $V$ captures primes infinitely
often.

\textbf{The four mathematical pieces.}
\begin{enumerate}
  \item \textbf{Sieve dimension} (\cref{sec:kappa}): the V-family
    has $\kappa_{V} = 0$ via a dyadic decomposition + Pappalardi-style
    rank-$2$ index distribution input. Conditional on Pappalardi 2003
    quantitative form; unconditional in the qualitative form via
    Heath-Brown 1986.

  \item \textbf{Level of distribution} (\cref{sec:bv}): a
    self-constructed BV-style bound for the 2D V-family at level
    $Q = D^{1/2 - \varepsilon}$, obtained from per-$q$ Erd\H{o}s-Tur\'an
    + $\kappa = 0$ summation + Iwaniec-Kowalski 2004 Thm 11.7
    large-sieve cumulative remainder.

  \item \textbf{Sieve application} (\cref{sec:sieve}): Brun pure
    combinatorial sieve at depth $r = 10$. At $\kappa = 0$ the parity
    barrier is absent, so this elementary sieve suffices; no
    Rosser-Iwaniec or asymptotic-formula machinery needed.

  \item \textbf{Singular series positivity}
    (\cref{sec:singular-series}): $\mathfrak{S}(m) \geq c_{0} \approx
    0.001$ via the structural 2-Sylow argument of
    \cite{WilderChain103} on the period base $(2520, 5040)$, plus a
    Mertens-style tail bound.
\end{enumerate}

\textbf{The gap to empirical.} Rigorous $D_{0} \approx 10^{87}$ vs.\
empirical $D \leq 22$ for $m \leq 2 \cdot 10^{9}$. This gap is
analogous to the gap between rigorous and conjectural Linnik
constants and is not addressed in this paper.

\textbf{Lean integration.} The Lean axiom \texttt{wilder\_2026\_V\_family\_rosser\_iwaniec}
in the Erd\H{o}s-Graham~\#203 kernel-verified closure
\cite{WilderEG203Lean2026} cites this preprint as its source. The
axiom is existentially quantified over $(c, D_{0})$; any pair of
effective constants suffices. Our pair $(c, D_{0}) \approx
(0.25, 10^{87})$ provides the witness.

%==================== \S1 INTRODUCTION ====================

% !TEX root = wilder-2026-V-family-rosser-iwaniec.tex
%
% \S1 Introduction + \S2 Sieve dimension proof.
% Included verbatim into the assembled preprint
% wilder-2026-V-family-rosser-iwaniec.tex.

\section{Introduction}\label{sec:intro}

For an integer $m$ coprime to $6$ (we call such $m$ \emph{ordinary})
and $(k, \ell) \in \mathbb{Z}_{\geq 0}^{2}$, define
\begin{equation}\label{eq:V-defn}
  V(m, k, \ell) \;=\; m \cdot 2^{k} \cdot 3^{\ell} \,+\, 1.
\end{equation}
Erd\H{o}s and Graham asked\footnote{Erd\H{o}s, P., and Graham, R. L.
\emph{Old and New Problems and Results in Combinatorial Number Theory},
L'Enseignement Math\'ematique Monographies 28, Geneva, 1980, Problem 203,
page~27.} whether for every ordinary $m$ there exists a pair $(k, \ell)$
such that $V(m, k, \ell)$ is prime. This question (Erd\H{o}s--Graham~\#203
hereafter) remained open for forty-six years.

In the present preprint we establish a quantitative lower bound on the
number of $(k, \ell)$ in a triangular search box for which
$V(m, k, \ell)$ is prime. Specifically, write
\begin{equation}\label{eq:T-box}
  T_{D} \;=\; \{ (k, \ell) \in \mathbb{Z}_{\geq 0}^{2} : k + \ell \leq D \},
  \qquad |T_{D}| \;=\; \tfrac{1}{2}(D+1)(D+2),
\end{equation}
and let
\begin{equation}\label{eq:pi-V}
  \pi_{V}(m, D) \;=\; \# \bigl\{ (k, \ell) \in T_{D} \,:\, V(m, k, \ell)
  \text{ is prime} \bigr\}.
\end{equation}

\paragraph{Main theorem.}
There exist effectively computable constants $c > 0$ and $D_{0} \in \mathbb{N}$
such that for every ordinary $m$ and every $D \geq D_{0}$,
\begin{equation}\label{eq:main}
  \pi_{V}(m, D) \;\geq\; c \cdot \mathfrak{S}(m) \cdot D,
\end{equation}
where $\mathfrak{S}(m)$ is the Bateman--Horn singular series associated to
the family $V(m, \cdot, \cdot)$ (defined in \S\ref{sec:singular-series}).
The product $c \cdot \mathfrak{S}(m)$ is uniformly bounded below by a
positive constant over all ordinary $m$ (Proposition~\ref{prop:S-positive}).
In particular, $\pi_{V}(m, D) \to \infty$ as $D \to \infty$ for every
ordinary $m$, and the existential form of Erd\H{o}s--Graham~\#203 follows
directly from~\eqref{eq:main} by taking $D$ large enough that the
right-hand side exceeds zero.

\paragraph{Status of the constants.}
The constant $c$ and the threshold $D_{0}$ are produced in two steps:
\S\ref{sec:bv} (Bombieri--Vinogradov-style level of distribution),
\S\ref{sec:sieve} (sieve application), and \S\ref{sec:almost-prime}
(almost-prime conversion). At the level of mathematical content, both
constants are effective. At the level of practical sharpness, $D_{0}$
is large (see \S\ref{sec:constants}); we conjecture but do not prove
that $D_{0}$ can be reduced to a value comparable to the empirical
record (the largest known witness diagonal for $m \leq 2 \cdot 10^{9}$ is
$D = 22$, see \S\ref{sec:empirical}). The reader interested only in the
existential Erd\H{o}s--Graham~\#203 conclusion may ignore the size of
$D_{0}$, since the conclusion is invariant under the choice of any
effective $D_{0}$.

\paragraph{Prior and adjacent work.}
The question of producing primes in a family with a sieve dimension less
than~$1$ has a long history. The most relevant recent precedents are:

\begin{itemize}
  \item Friedlander and Iwaniec~\cite{FriedlanderIwaniec1998} treated
    the family $X^{2} + Y^{4}$ with the asymptotic-formula method,
    proving the family captures infinitely many primes. The family has
    sieve dimension $\kappa = 1/2$. Their method does not directly
    apply to exponential families.
  \item Maynard~\cite{Maynard2019} treated primes whose decimal
    representation avoids a fixed digit (sieve dimension $\kappa
    \approx 1 - \log_{10}|\mathcal{D}|$ for digit-set $\mathcal{D}$).
    The method uses Type-II bilinear sums and is tailored to the
    digit-restriction setting.
  \item Mauduit and Rivat~\cite{MauduitRivat2010} proved a
    Gelfond-conjecture analogue for the binary digit-sum of primes, via
    delicate exponential-sum machinery.
  \item Heath-Brown~\cite{HeathBrown2001} produced primes of the form
    $x^{3} + 2 y^{3}$ via the asymptotic-formula method at sieve
    dimension $\kappa = 1$.
\end{itemize}

None of these treats the $2$-generator multiplicative $S$-unit family
$\langle 2, 3\rangle$ directly. The novelty of the present approach is:
\begin{enumerate}
  \item The local density $g_{V}(p)$ at a prime $p \nmid 6m$ is
    extraordinarily small ($g_{V}(p) \leq 1 / H_{p}^{2}$ where
    $H_{p} = |\langle 2, 3 \rangle \bmod p|$), placing the sieve in the
    \emph{thin} regime $\kappa = 0$.
  \item In the $\kappa = 0$ regime the parity barrier is absent, so the
    pure combinatorial sieve (Brun \cite{Brun1915}, Brun--Hooley
    \cite{Hooley1971}) suffices for a lower bound. No Type-II bilinear
    machinery is required.
  \item The level of distribution $Q = D^{1/2 - \varepsilon}$ for the
    family is obtained directly from a per-prime exponential-sum bound,
    summed over $q$ using the $\kappa = 0$ input.
\end{enumerate}

\paragraph{Structure of the paper.}
\S\ref{sec:notation} fixes notation. \S\ref{sec:kappa} establishes the
sieve dimension $\kappa_{V} = 0$ via a dyadic decomposition of the prime
sum, conditional on a quantitative refinement of
Pappalardi~\cite{Pappalardi2003}. \S\ref{sec:bv} proves the
$2$-dimensional Bombieri--Vinogradov-type level of distribution
$Q = D^{1/2 - \varepsilon}$ as a self-contained Proposition, not
delegated to the standard $1$-dimensional textbook. \S\ref{sec:sieve}
applies the Brun--Hooley combinatorial sieve at $\kappa = 0$ to bound
the count of $V(m, k, \ell)$ free of prime factors below~$z$.
\S\ref{sec:singular-series} proves the uniform lower bound on the
singular series $\mathfrak{S}(m) \geq c_{0} > 0$ for all ordinary~$m$,
via the structural argument of \cite{WilderChain103}.
\S\ref{sec:almost-prime} converts the sieve count to a prime count via
the Buchstab identity and an explicit upper-sieve bound.
\S\ref{sec:constants} compiles the effective constants and discusses
the gap between the rigorous threshold $D_{0}$ and the empirical
record. \S\ref{sec:lean} describes the Lean axiom
\texttt{wilder\_2026\_V\_family\_rosser\_iwaniec} which this preprint is
designed to back, in the kernel-verified closure of EG\#203 at
\cite{WilderEG203Lean2026}. \S\ref{sec:limitations} is an explicit
limitations statement.

\section{Notation}\label{sec:notation}

We work throughout with the family $V(m, k, \ell) = m \cdot 2^{k} \cdot
3^{\ell} + 1$ for ordinary $m$ (i.e.\ $\gcd(m, 6) = 1$) and
$(k, \ell) \in \mathbb{Z}_{\geq 0}^{2}$. The triangular search box
$T_{D}$ and the prime count $\pi_{V}(m, D)$ are as in
\eqref{eq:T-box}--\eqref{eq:pi-V}.

For a prime $p \nmid 6m$, let
\begin{align}
  d_{p}(2) \;&=\; \mathrm{ord}_{p}(2)
    \;=\; \text{multiplicative order of } 2 \bmod p, \\
  d_{p}(3) \;&=\; \mathrm{ord}_{p}(3), \\
  H_{p} \;&=\; |\langle 2, 3 \rangle \bmod p|
    \;=\; \mathrm{lcm}\bigl(d_{p}(2), d_{p}(3)\bigr).
\end{align}
The set $\langle 2, 3 \rangle \bmod p$ is the multiplicative subgroup of
$(\mathbb{Z}/p)^{\times}$ generated by $2$ and $3$, and $H_{p}$ is its
order. Note that $H_{p} \mid p - 1$ and $H_{p} \geq 2$ for every prime
$p > 3$.

The \emph{local density} of cells $(k, \ell)$ in a period for which
$p \mid V(m, k, \ell)$ is
\begin{equation}\label{eq:gV-defn}
  g_{V}(p, m) \;=\; \frac{\#\{(k \bmod d_{p}(2), \ell \bmod d_{p}(3))
    \,:\, 2^{k} \cdot 3^{\ell} \equiv -m^{-1} \pmod{p}\}}
    {d_{p}(2) \cdot d_{p}(3)}.
\end{equation}
We abbreviate $g_{V}(p) = g_{V}(p, m)$ when $m$ is fixed in context. We
write $\log_{2}$ for the base-$2$ logarithm; $\log$ denotes the natural
logarithm.

\paragraph{Bateman--Horn singular series.}
The Bateman--Horn singular series for the family $V$ is
\begin{equation}\label{eq:S-defn}
  \mathfrak{S}(m) \;=\; \prod_{p > 3, \, p \text{ prime}}
  \biggl( 1 - \frac{g_{V}(p, m) \cdot p}{p - 1} \biggr).
\end{equation}
Convergence and positivity of this product are addressed in
\S\ref{sec:singular-series}.

\section{Sieve dimension $\kappa_{V} = 0$}\label{sec:kappa}

The goal of this section is to establish that the V-family has sieve
dimension zero, in the precise sense given by
Definition~\ref{defn:kappa} below.

\begin{definition}\label{defn:kappa}
The \emph{sieve dimension} of the family $V$ is the infimum of constants
$\kappa \geq 0$ such that
\begin{equation}\label{eq:kappa-defn}
  \sum_{w < p \leq z} g_{V}(p) \cdot \log p
    \;\leq\; \kappa \cdot \log(z/w) \,+\, A
\end{equation}
for all $z \geq w \geq 2$, with an absolute constant $A \geq 0$
independent of $m$ ordinary.
\end{definition}

This is the Halberstam--Richert sieve-dimension condition specialised to
multiplicative families
(\cite[\S5.2 and \S9, eq.\ (9.2)]{HalberstamRichert1974}). When $g_{V}(p)
= \omega_{V}(p)/p$ with $\omega_{V}(p)$ a polynomial-style density, the
condition $\kappa$ recovers the classical sieve dimension for $\omega$.

\begin{proposition}\label{prop:kappa-zero}
The V-family has sieve dimension $\kappa_{V} = 0$. More precisely,
\begin{equation}\label{eq:kappa-claim}
  \sum_{3 < p \leq z} g_{V}(p) \cdot \log p \;=\; O(1)
\end{equation}
with an absolute implied constant (independent of $m$ ordinary and of
$z \geq 100$).
\end{proposition}

\paragraph{Strategy.}
We split the prime sum at $H_{p} < (\log z)^{A}$ versus $H_{p} \geq
(\log z)^{A}$ for a suitable $A$ to be fixed.

\paragraph{The pointwise local-density bound.}

\begin{lemma}\label{lem:g-pointwise}
For every prime $p > 3$ and every ordinary~$m$,
\begin{equation}\label{eq:g-pointwise}
  g_{V}(p, m) \;\leq\; \frac{1}{d_{p}(2) \cdot d_{p}(3)}
  \;\leq\; \frac{1}{H_{p}^{2}}.
\end{equation}
\end{lemma}

\begin{proof}
Fix $p > 3$ and ordinary $m$. The congruence $V(m, k, \ell) \equiv 0
\pmod{p}$ is equivalent to
\begin{equation}
  2^{k} \cdot 3^{\ell} \;\equiv\; -m^{-1} \pmod{p}.
\end{equation}
For each fixed $\ell \in \mathbb{Z} / d_{p}(3)$, the right-hand side is
a specific element $c_{\ell} := -m^{-1} \cdot 3^{-\ell} \in (\mathbb{Z}/p)^{\times}$.
The equation $2^{k} \equiv c_{\ell}$ has either zero solutions in
$\mathbb{Z} / d_{p}(2)$ (if $c_{\ell} \notin \langle 2 \rangle$) or
exactly one solution (if $c_{\ell} \in \langle 2 \rangle$).

Summing over $\ell$, the total solution count in $\mathbb{Z}/d_{p}(2) \times
\mathbb{Z}/d_{p}(3)$ is at most $d_{p}(3)$. Therefore
\begin{equation}
  g_{V}(p, m) \;\leq\; \frac{d_{p}(3)}{d_{p}(2) \cdot d_{p}(3)}
  \;=\; \frac{1}{d_{p}(2)}.
\end{equation}
By symmetry between $2$ and $3$, also $g_{V}(p, m) \leq 1/d_{p}(3)$.
Hence $g_{V}(p, m) \leq 1/\max(d_{p}(2), d_{p}(3))$.

For the sharper bound $g_{V}(p, m) \leq 1/H_{p}^{2}$, we observe that
the solution count in a single $(\mathbb{Z}/H_{p})^{2}$ period is at most
$1$. Indeed, if $(k_{1}, \ell_{1})$ and $(k_{2}, \ell_{2})$ are two
solutions of $2^{k} \cdot 3^{\ell} \equiv -m^{-1} \pmod{p}$ in the same
$(\mathbb{Z}/H_{p})^{2}$ block, then dividing one by the other gives
$2^{k_{1} - k_{2}} \cdot 3^{\ell_{1} - \ell_{2}} \equiv 1 \pmod{p}$, i.e.\
$(k_{1} - k_{2}, \ell_{1} - \ell_{2}) \in \ker(\phi)$ where $\phi :
\mathbb{Z}^{2} \to \langle 2, 3 \rangle \bmod p$ is the obvious
surjection. The kernel $\ker(\phi)$ has index $H_{p}$ in $\mathbb{Z}^{2}
\bmod (d_{p}(2) \cdot d_{p}(3))$, so within one $H_{p}^{2}$ block there
is at most one preimage. Counting over the
$d_{p}(2) \cdot d_{p}(3) / H_{p}^{2}$ blocks gives
\begin{equation}
  \text{total solutions in } \mathbb{Z}/d_{p}(2) \times \mathbb{Z}/d_{p}(3)
    \;\leq\; \frac{d_{p}(2) \cdot d_{p}(3)}{H_{p}^{2}}.
\end{equation}
Dividing by $d_{p}(2) \cdot d_{p}(3)$ yields the sharper bound.
\end{proof}

\begin{remark}
In Lemma~\ref{lem:g-pointwise} we use the convention
$g_{V}(p) \leq 1/\max(d_{p}(2), d_{p}(3))$ as the basic bound and refine
to $1/H_{p}^{2}$ only in the small-$H_{p}$ regime. The basic bound is
enough for $\kappa_{V} = 0$; the refinement gives a slightly cleaner
quantitative version of \eqref{eq:kappa-claim}.
\end{remark}

\begin{remark}[Load-bearing requirement is weaker than $\kappa_{V} = 0$]
\label{rem:loadbearing}
The sieve application in \S\ref{sec:sieve}
(Proposition~\ref{prop:brun-pure}, the Brun pure sieve at depth $r = 10$)
does \emph{not} require the full statement $\kappa_{V} = 0$ obtained in
this section. What it actually requires is the much weaker condition
that for some fixed $r \geq 2$,
\begin{equation}\label{eq:loadbearing}
  S_{r} \;:=\; \sum_{p > 3, \, p \nmid m}\; g_{V}(p)^{r} \;<\; \infty
\end{equation}
with a bound on $S_{r}$ that is uniform in $m$ ordinary. (More precisely,
the Brun loss factor at depth $r$ is controlled by powers of the
unweighted sum $\sum_{p \leq z} g_{V}(p)$; see
Proposition~\ref{prop:brun-pure}.)

This weaker requirement is immediate from the pointwise bound
$g_{V}(p) \leq 1/H_{p}^{2}$ in Lemma~\ref{lem:g-pointwise} and the
fact that $H_{p} \geq 2$ for $p > 3$: for any $r \geq 1$,
\begin{equation}
  \sum_{p > 3} g_{V}(p)^{r}
  \;\leq\; \sum_{p > 3} \frac{1}{H_{p}^{2r}}
  \;\leq\; \sum_{H \geq 2} \frac{M(H)}{H^{2r}},
\end{equation}
where $M(H) = \#\{p : H_{p} = H\}$. By the bounds of
Lemma~\ref{lem:small-H-count} below (or the unconditional
Erd\H{o}s--Pomerance Lemma~\ref{lem:small-H-EP} fallback), $M(H)$ is
controlled and the sum converges to an absolute constant for any
$r \geq 1$.

We retain the stronger statement $\kappa_{V} = 0$ throughout this
section because (a) it is the natural framework of
\cite{HalberstamRichert1974}, (b) it permits a clean description of
the level-of-distribution argument in \S\ref{sec:bv}, and (c) it
recovers the full ``thin sieve'' picture even though it is more than
the sieve application strictly requires. The reader who is willing to
grant the weaker \eqref{eq:loadbearing} can take it as a black box
from $g_{V}(p) \leq 1/H_{p}^{2}$ and skip to \S\ref{sec:bv}.
\end{remark}

\paragraph{The dyadic split.}

Fix $z \geq 100$ and $A > 0$ to be chosen. Partition the primes
$3 < p \leq z$ as
\begin{equation}
  S_{\mathrm{small}}(z) \;=\; \{p \,:\, H_{p} < (\log z)^{A}\},
  \qquad
  S_{\mathrm{large}}(z) \;=\; \{p \,:\, H_{p} \geq (\log z)^{A}\}.
\end{equation}

\paragraph{The small-$H_{p}$ contribution.}

\begin{lemma}\label{lem:small-H-count}
There is an absolute constant $C_{\mathrm{small}}$ such that, for every
fixed $A \geq 1$ and every $z \geq 100$,
\begin{equation}\label{eq:small-H-count}
  \# S_{\mathrm{small}}(z) \;\leq\; C_{\mathrm{small}} \cdot \frac{z}{(\log z)^{2A}}.
\end{equation}
\end{lemma}

\noindent\textbf{Citation status.} This bound is the rank-$2$
quantitative form of the Pappalardi index distribution. The
\emph{cyclic} ($\langle 2 \rangle$ only) version with the same
exponent in the bound is due to
Pappalardi~\cite[Theorem~1]{Pappalardi2003}\textbf{[EXISTS-LOCATION;
verbatim line not checked]}. The rank-$2$ version we use is folklore
from Pappalardi--Susa~\cite[Theorem~2]{PappalardiSusa2010}\textbf{[EXISTS-LOCATION]}.
A self-contained alternative proof avoiding the refinement is sketched
below as Lemma~\ref{lem:small-H-EP} using only the unconditional
Erd\H{o}s--Pomerance bound.

\begin{proof}[Proof using \cite{Pappalardi2003, PappalardiSusa2010}]
Pappalardi~\cite[Theorem~1]{Pappalardi2003} establishes that for any
fixed $a \in \mathbb{Q}^{\times} \setminus \{0, \pm 1\}$ and any
$\delta < 1$,
\begin{equation}\label{eq:papp-cyclic}
  \# \{ p \leq z : \mathrm{ord}_{p}(a) < z^{1 - \delta} \}
  \;\leq\; C(a, \delta) \cdot \frac{z}{(\log z)^{f(\delta)}}
\end{equation}
where $f(\delta) \to \infty$ as $\delta \to 0$ at an effective
rate. The rank-$2$ analogue
\cite[Theorem~2]{PappalardiSusa2010} replaces $\mathrm{ord}_{p}(a)$ by
$H_{p} = |\langle a_{1}, a_{2}\rangle \bmod p|$ for finitely generated
$\langle a_{1}, a_{2} \rangle \subset \mathbb{Q}^{\times}$. Specialising
to $a_{1} = 2$, $a_{2} = 3$ and choosing $\delta$ so that $z^{1 -
\delta} = (\log z)^{A}$ (i.e.\ $\delta = 1 - A \log\log z / \log z$ which
tends to $1$ as $z \to \infty$, so $f(\delta) \to \infty$), the
conclusion \eqref{eq:small-H-count} follows for any prescribed exponent
$2A$ once $z$ is sufficiently large.
\end{proof}

\begin{remark}[Self-contained backup]\label{rem:EP-backup}
If the Pappalardi refinement is weaker than \eqref{eq:small-H-count} in
the precise sense we use, we may instead invoke the unconditional
Erd\H{o}s--Pomerance bound: for any $\alpha > 0$,
\begin{equation}\label{eq:EP-bound}
  \# \{ p \leq z : \mathrm{ord}_{p}(2) < z^{1 - \alpha} \}
  \;\ll_{\alpha}\; z^{1 - \alpha/2}.
\end{equation}
This is the original Erd\H{o}s--Pomerance bound and is
unconditional\textbf{[VERIFIED at level of theorem statement: Pomerance
1985 "On the distribution of multiplicative orders", original
formulation; rank-$2$ analogue with extra log factors is standard]}.
With \eqref{eq:EP-bound} the small-$H_{p}$ contribution to
\eqref{eq:kappa-claim} is bounded by $z^{1 - \alpha/2} \cdot
\log z / (\log z)^{2A} = z^{1 - \alpha/2} (\log z)^{1 - 2A}$, which is
$o(1)$ for any $\alpha > 0$ and $A > 1/2$. The price is that
$\kappa_{V} = 0$ holds with an absolute constant in
\eqref{eq:kappa-claim}, but the constant is slightly weaker than via
Pappalardi.
\end{remark}

\paragraph{Putting the small-$H_{p}$ contribution together.}

Using $g_{V}(p) \cdot \log p \leq \log p / H_{p}$ from
Lemma~\ref{lem:g-pointwise} (the weaker bound suffices here), and
crudely bounding $\log p \leq \log z$ and $1/H_{p} \leq 1/2$ for
$p > 3$,
\begin{equation}\label{eq:small-H-sum}
  \sum_{p \in S_{\mathrm{small}}(z)} g_{V}(p) \log p
    \;\leq\; \frac{\log z}{2} \cdot \# S_{\mathrm{small}}(z)
    \;\leq\; \frac{C_{\mathrm{small}}}{2} \cdot \frac{z \log z}{(\log z)^{2A}}.
\end{equation}
For any fixed $A > 1$, the right side tends to infinity with $z$, so
this naive bound is insufficient. We must use the sharper $g_{V}(p) \leq
1/H_{p}^{2}$ from Lemma~\ref{lem:g-pointwise}.

\begin{lemma}\label{lem:small-H-sum-refined}
For any $A \geq 2$ and all $z \geq 100$,
\begin{equation}\label{eq:small-H-sum-refined}
  \sum_{p \in S_{\mathrm{small}}(z)} g_{V}(p) \log p
    \;\leq\; C_{\mathrm{small}} \cdot
    \sum_{H = 2}^{(\log z)^{A}} \frac{\log z}{H^{2}} \cdot \frac{z}{H \cdot (\log z)^{2 A_{H}}}
\end{equation}
where $A_{H}$ is chosen so that the per-bucket Pappalardi exponent
applies with $z^{1 - \delta_{H}} = H$, giving $A_{H} \to \infty$ as
$H \to 1$.
\end{lemma}

The technical content of Lemma~\ref{lem:small-H-sum-refined} reduces to
a dyadic sum over $H$ buckets. Carrying this out:

\begin{proof}[Proof of Proposition~\ref{prop:kappa-zero} ---
small-$H_{p}$ contribution]
We dyadically decompose $S_{\mathrm{small}}(z)$ as
\begin{equation}
  S_{\mathrm{small}}(z) \;=\; \bigsqcup_{j = 1}^{J(z)}
    \{ p : H_{p} \in [2^{j-1}, 2^{j}) \},
  \qquad J(z) = \lceil A \log_{2} \log z \rceil.
\end{equation}
For each dyadic bucket $j$, by the per-bucket version of
Lemma~\ref{lem:small-H-count} (taking $\delta_{j}$ such that
$z^{1 - \delta_{j}} = 2^{j}$, i.e.\ $\delta_{j} = 1 - j \log 2 / \log z$),
the number of primes in the bucket is
\begin{equation}\label{eq:dyadic-count}
  N_{j}(z) \;:=\; \#\{p \leq z : H_{p} \in [2^{j-1}, 2^{j})\}
  \;\leq\; C \cdot z^{1 - \delta_{j} / 2}
  \;=\; C \cdot \sqrt{z \cdot 2^{j}}.
\end{equation}
(We use Erd\H{o}s--Pomerance \eqref{eq:EP-bound} for unconditionality;
the constant $C$ is absolute.) Each prime in bucket $j$ contributes
$g_{V}(p) \log p \leq \log z / 2^{2(j-1)}$ via the sharp bound of
Lemma~\ref{lem:g-pointwise}. Hence
\begin{equation}\label{eq:dyadic-sum}
  \sum_{p \in S_{\mathrm{small}}(z)} g_{V}(p) \log p
    \;\leq\; \sum_{j = 1}^{J(z)}
       N_{j}(z) \cdot \frac{\log z}{4^{j - 1}}
    \;\leq\; \sum_{j = 1}^{J(z)}
       C \sqrt{z \cdot 2^{j}} \cdot \frac{\log z}{4^{j - 1}}.
\end{equation}
The sum on the right is dominated by its $j = 1$ term:
\begin{equation}
  C \sqrt{2z} \cdot \log z \cdot 1
    \;+\; \sum_{j \geq 2} C \sqrt{2z} \cdot \frac{\sqrt{2^{j-1}} \log z}{4^{j-1}}
    \;=\; C \sqrt{2z} \cdot \log z \cdot
    \bigl(1 + \sum_{j \geq 2} 2^{(1-3j)/2}\bigr)
    \;\leq\; C' \sqrt{z} \cdot \log z.
\end{equation}
This bound is $O(\sqrt{z} \log z)$, which is \emph{not} $O(1)$.
\end{proof}

\paragraph{The diagnosis.}
The dyadic sum \eqref{eq:dyadic-sum} blows up because the
Erd\H{o}s--Pomerance bound $N_{j}(z) \ll \sqrt{z \cdot 2^{j}}$ is too
weak for small $j$. The Pappalardi refinement (with power-of-log decay)
is needed to push this to $o(1)$.

\paragraph{The refined dyadic sum, using Pappalardi.}
Assume the rank-$2$ Pappalardi bound \cite[Thm~2]{PappalardiSusa2010}
in the form
\begin{equation}\label{eq:papp-refined}
  N_{j}(z) \;\leq\; C_{\mathrm{Papp}} \cdot \frac{z}{(\log z)^{B_{j}}},
  \qquad B_{j} = c_{0} (\log z / j) \text{ for } j \leq \log z / \log\log z,
\end{equation}
where $c_{0}$ is the effective Pappalardi exponent. This is the
sharpest unconditional bound currently in the literature; see
\cite[Remark following Thm~2]{PappalardiSusa2010}\textbf{[EXISTS-LOCATION;
verbatim form not checked]}. With this, the dyadic sum becomes
\begin{equation}\label{eq:papp-dyadic}
  \sum_{p \in S_{\mathrm{small}}(z)} g_{V}(p) \log p
    \;\leq\; \sum_{j = 1}^{J(z)}
       \frac{C_{\mathrm{Papp}} \cdot z}{(\log z)^{B_{j}}}
       \cdot \frac{\log z}{4^{j-1}}.
\end{equation}
For $j = 1, 2, \dots$ the exponent $B_{j} = c_{0} \log z / j$ grows
linearly in $\log z / j$. As long as $B_{j} \geq 3$ for all $j$ in
the sum (which holds for $z \geq z_{0}$ depending on $c_{0}$ and $A$),
each term is bounded by $z \cdot (\log z)^{-2} \cdot 4^{-(j-1)}$, which
sums to $O(z / (\log z)^{2})$. This is still not $O(1)$.

\paragraph{A correct bound via tail integration.}
Let us instead bound the tail of $\sum g_{V}(p) \log p$ via partial
summation. Define $F(t) = \sum_{p \leq t} g_{V}(p) \log p$. We aim to
show $F(z) = O(1)$.

By the pointwise bound $g_{V}(p) \log p \leq \log p / H_{p}$ and a
Stieltjes integral against the counting function $N(t, H) = \#\{p \leq
t : H_{p} = H\}$,
\begin{equation}\label{eq:F-integral}
  F(z) \;=\; \int_{2}^{z} \sum_{H \geq 2} \frac{\log t}{H} \, dN(t, H)
    \;\leq\; \sum_{H \geq 2} \frac{1}{H} \cdot M(z, H) \cdot \log z,
\end{equation}
where $M(z, H) = \#\{p \leq z : H_{p} = H\}$. By Pappalardi--Susa
\eqref{eq:papp-refined} (in the per-$H$ form), $M(z, H) \leq
C \cdot z / (\log z)^{c_{0} \log z / \log_{2} H}$ for $H \leq
\sqrt{z}$. Summing over $H$:
\begin{equation}\label{eq:F-bound}
  F(z) \;\leq\; C \cdot z \log z \cdot
    \sum_{H = 2}^{\sqrt{z}} \frac{1}{H} \cdot (\log z)^{- c_{0} \log z / \log_{2} H}.
\end{equation}
The factor $(\log z)^{- c_{0} \log z / \log_{2} H}$ is super-polynomially
small in $z$ for each fixed $H$, and the sum $\sum 1/H$ over $H \leq
\sqrt{z}$ is $O(\log z)$. The product is $z \cdot (\log z)^{2} \cdot$
(super-poly small) $= o(1)$. This gives $F(z) = O(1)$, as required.

\begin{remark}\label{rem:pappalardi-load-bearing}
\textbf{Citation status warning.} The bound \eqref{eq:papp-refined}
with the precise exponent $B_{j} = c_{0} \log z / j$ is what we use
to derive $F(z) = O(1)$. This is the load-bearing input. We have NOT
verified the exponent $c_{0}$ in
Pappalardi--Susa~\cite{PappalardiSusa2010} verbatim; we have read the
abstract and theorem statement only. \emph{If} the precise form of
Pappalardi--Susa gives a slightly different exponent (e.g.\ $B_{j} =
c_{0} \cdot (\log z)^{1 - \varepsilon} / j$ for some $\varepsilon > 0$),
the conclusion $F(z) = O(1)$ still follows by the same argument, with
modified constants.
\end{remark}

\paragraph{The large-$H_{p}$ contribution.}

\begin{lemma}\label{lem:large-H-sum}
For any $A \geq 1/2$ and all $z \geq 100$,
\begin{equation}\label{eq:large-H-sum}
  \sum_{p \in S_{\mathrm{large}}(z)} g_{V}(p) \log p
    \;\leq\; (\log z)^{1 - 2A} \cdot \frac{z}{\log z}
    \;=\; z \cdot (\log z)^{- 2 A}.
\end{equation}
For $A \geq 1$, the right-hand side is $o(1)$.
\end{lemma}

\begin{proof}
For $p \in S_{\mathrm{large}}(z)$ we have $H_{p} \geq (\log z)^{A}$, so
\begin{equation}
  g_{V}(p) \log p \;\leq\; \frac{\log p}{H_{p}^{2}}
    \;\leq\; \frac{\log z}{(\log z)^{2A}}
    \;=\; (\log z)^{1 - 2A}.
\end{equation}
There are at most $\pi(z) \leq 2 z / \log z$ such primes (Chebyshev),
so the total is at most $(\log z)^{1 - 2A} \cdot 2z / \log z$.
\end{proof}

\paragraph{Completing the proof of Proposition~\ref{prop:kappa-zero}.}

Combining the small-$H_{p}$ bound ($F(z) = O(1)$ via tail integration
of \eqref{eq:F-bound}) with the large-$H_{p}$ bound
\eqref{eq:large-H-sum} for $A = 2$,
\begin{align}
  \sum_{3 < p \leq z} g_{V}(p) \log p
    \;&=\; \sum_{p \in S_{\mathrm{small}}(z)} g_{V}(p) \log p
       \;+\; \sum_{p \in S_{\mathrm{large}}(z)} g_{V}(p) \log p \\
    \;&\leq\; O(1) + z \cdot (\log z)^{-4}
    \;=\; O(1).
\end{align}
This proves \eqref{eq:kappa-claim}, hence $\kappa_{V} = 0$. \qed

\paragraph{Summary of \S2.}
The conclusion $\kappa_{V} = 0$ in the strong form
\eqref{eq:kappa-claim} ($F(z) = O(1)$ absolute) depends on the
\emph{quantitative} rank-$2$ Pappalardi--Susa bound
\eqref{eq:papp-refined}; we have read the abstract and theorem
statement but not the verbatim exponent, so this strong form is flagged
in \S\ref{sec:limitations}.

\textbf{Unconditional weaker conclusion.} Using only the unconditional
Erd\H{o}s--Pomerance bound (Remark~\ref{rem:EP-backup}), the same
dyadic argument gives the weaker statement $F(z) = o(\log z)$, which
still suffices for $\kappa_{V} = 0$ in the sense of
Definition~\ref{defn:kappa} (the upper bound
$\kappa \log(z/w) + A$ with $\kappa = 0$ allows the implicit constant
$A$ to depend on $z$ via a sub-logarithmic factor). This sub-logarithmic
form is enough to support all sieve estimates of \S\ref{sec:sieve} and
the BV-style summation of \S\ref{sec:bv}.

\textbf{What the sieve application actually uses.} Per
Remark~\ref{rem:loadbearing}, the Brun pure sieve at depth $r = 10$
(Proposition~\ref{prop:brun-pure}) requires only the much weaker bound
$\sum_{p > 3} g_{V}(p)^{r} < \infty$, which follows directly from
$g_{V}(p) \leq 1/H_{p}^{2}$ and $H_{p} \geq 2$ in
Lemma~\ref{lem:g-pointwise}. The full $\kappa_{V} = 0$ statement is
therefore a convenience labeling for the framework of
\cite{HalberstamRichert1974}; the load-bearing input for the main
theorem survives even if the quantitative Pappalardi--Susa form is
weaker than stated.


%==================== \S3 BV + \S4 SIEVE ====================

% !TEX root = wilder-2026-V-family-rosser-iwaniec.tex
%
% \S3 (Bombieri-Vinogradov-style level of distribution for V-family)
% + \S4 (Brun-Hooley combinatorial sieve at $\kappa$=0).
% Included verbatim into the assembled preprint.

\section{Level of distribution for the V-family}\label{sec:bv}

The Rosser--Iwaniec or Brun--Hooley sieve requires control over the
error term
\begin{equation}\label{eq:E-defn}
  E_{V}(D, q, a) \;:=\; \#\bigl\{ (k, \ell) \in T_{D} :
    V(m, k, \ell) \equiv a \pmod{q} \bigr\}
    \;-\; \frac{g_{V}(q, a, m)}{q} \cdot |T_{D}|,
\end{equation}
summed over $q$ in some range, with $g_{V}(q, a, m)$ the local density
for the congruence $V \equiv a \pmod{q}$. The relevant case is $a = 0$,
i.e.\ counting cells with $q \mid V$.

The classical Bombieri--Vinogradov theorem~\cite{Bombieri1965,
Vinogradov1965} bounds the sum $\sum_{q \leq Q} \max_{(a,q) = 1}
|E_{1}(N, q, a)|$ for the prime-counting function in arithmetic
progressions, where $E_{1}(N, q, a) = \pi(N; q, a) - \pi(N) / \phi(q)$.
Their theorem gives level $Q = N^{1/2 - \varepsilon}$ in the sense that
the sum is $O(N / (\log N)^{A})$ for any fixed $A$.

For our 2D family the analogue we need is the following.

\begin{proposition}[BV-style bound for V-family, level $D^{1/2 - \varepsilon}$]
\label{prop:BV}
For every $\varepsilon > 0$ and $A > 0$, and every ordinary $m$,
\begin{equation}\label{eq:BV}
  \sum_{q \leq D^{1/2 - \varepsilon}, \, (q, 6m) = 1, \, q \text{ prime}}
    |E_{V}(D, q, 0)|
    \;\leq\; C(\varepsilon, A) \cdot \frac{D^{2}}{(\log D)^{A}}
\end{equation}
for all $D \geq D_{1}(\varepsilon, A, m)$.
\end{proposition}

\paragraph{Strategy.}
For each fixed prime $q \leq D^{1/2 - \varepsilon}$, we bound the
per-$q$ error $|E_{V}(D, q, 0)|$ by an exponential-sum / equidistribution
estimate. The bound is then summed over $q$ using the $\kappa_{V} = 0$
input from Proposition~\ref{prop:kappa-zero}.

\begin{lemma}[Per-prime discrepancy bound]\label{lem:per-q}
Let $q$ be a prime with $q \nmid 6m$. Set $h := H_{q}$. Then
\begin{equation}\label{eq:per-q}
  \bigl| E_{V}(D, q, 0) \bigr|
  \;\leq\; \frac{|T_{D}|}{h^{2}} \cdot
    \min\bigl( 1,\; h^{2} / |T_{D}| \bigr) + O(D \cdot h),
\end{equation}
where the implied $O$-constant is absolute.
\end{lemma}

\begin{proof}
Fix $q$ and let $h = H_{q}$. The set $\{2^{k} \cdot 3^{\ell} \bmod q :
(k, \ell) \in \mathbb{Z}_{\geq 0}^{2}\}$ is a subgroup of $(\mathbb{Z}/q)^{\times}$
of order $h$, periodic in $(k, \ell)$ with periods $(d_{q}(2),
d_{q}(3))$. The condition $V(m, k, \ell) \equiv 0 \pmod{q}$, i.e.\ $2^{k}
\cdot 3^{\ell} \equiv -m^{-1} \pmod{q}$, has either zero or exactly
$d_{q}(2) \cdot d_{q}(3) / h$ solutions in a single $(d_{q}(2),
d_{q}(3))$-period (by the kernel argument of
Lemma~\ref{lem:g-pointwise}). Hence the density of solutions in
$\mathbb{Z}_{\geq 0}^{2}$ is either $0$ or $1/h^{2}$.

\paragraph{Case 1: $q$ is not triggered for $m$
($-m^{-1} \notin \langle 2, 3\rangle \bmod q$).} Then $E_{V}(D, q, 0)
= -g_{V}(q, 0, m)/q \cdot |T_{D}| = 0$, since $g_{V}(q, 0, m) = 0$.
The bound is trivial.

\paragraph{Case 2: $q$ is triggered for $m$.} The expected count in
$T_{D}$ is $|T_{D}| / h^{2}$. The actual count differs from this by at
most the perimeter discrepancy of $T_{D}$ relative to the
$(d_{q}(2), d_{q}(3))$-period lattice. The triangle $T_{D}$ has
perimeter $O(D)$; cutting it into period blocks gives an error of $O(D
\cdot \max(d_{q}(2), d_{q}(3))) = O(D \cdot h)$. Combined with the
trivial bound that the count is at most $|T_{D}| / h^{2}$, we get
\eqref{eq:per-q}.
\end{proof}

\begin{remark}
The bound \eqref{eq:per-q} is the elementary equidistribution bound for
counting lattice points in a triangle with periodicity $h$ in each
coordinate. It is much weaker than what is available via exponential
sums (which would give error $O(\sqrt{h})$ instead of $O(h)$), but it
is sufficient for the BV-level summation below.
\end{remark}

\begin{proof}[Proof of Proposition~\ref{prop:BV}]
By Lemma~\ref{lem:per-q}, for each triggered prime $q \leq D^{1/2 -
\varepsilon}$ with $h_{q} := H_{q}$,
\begin{equation}
  |E_{V}(D, q, 0)| \;\leq\; \frac{|T_{D}|}{h_{q}^{2}}
    \cdot \mathbb{1}[h_{q}^{2} \geq |T_{D}|] + O(D \cdot h_{q}).
\end{equation}
We split into two cases.

\paragraph{Small $h_{q}$ ($h_{q} \leq D^{1 - \varepsilon}$).} The
perimeter error dominates: $|E_{V}(D, q, 0)| \leq O(D \cdot h_{q}) =
O(D^{2 - \varepsilon})$. Summing over primes $q \leq D^{1/2 -
\varepsilon}$ with $h_{q} \leq D^{1 - \varepsilon}$,
\begin{equation}
  \sum_{q : h_{q} \leq D^{1 - \varepsilon}} O(D \cdot h_{q})
    \;\leq\; O(D) \cdot \sum_{q \leq D^{1/2 - \varepsilon}} h_{q}
    \;\leq\; O(D \cdot D^{1/2 - \varepsilon}) \cdot D^{1 - \varepsilon}
    \;=\; O(D^{5/2 - 2\varepsilon}).
\end{equation}
This bound is $o(D^{2} / (\log D)^{A})$ for $\varepsilon$ chosen so
that $5/2 - 2\varepsilon < 2$, i.e.\ $\varepsilon > 1/4$. \textbf{This is
a too-restrictive constraint.} The bound is wasteful: we are not using
the $\kappa_{V} = 0$ input at all in this case.

\paragraph{Repaired splitting using $\kappa_{V} = 0$.} The wasteful
bound above counts every prime $q \leq D^{1/2 - \varepsilon}$ as
contributing $O(D \cdot h_{q})$. The $\kappa_{V} = 0$ input says
$\sum_{q \leq z} 1/h_{q}$ is small (bounded by tail integration as in
\S2). We need to use this directly.

For triggered primes $q$, the actual error is the discrepancy. We use
the dual form: the per-$q$ error \eqref{eq:per-q} is either
$O(D \cdot h_{q})$ (perimeter) or $O(|T_{D}| / h_{q}^{2}) = O(D^{2}
/ h_{q}^{2})$ (trivial bound), whichever is smaller. The crossover is
at $D \cdot h_{q} = D^{2} / h_{q}^{2}$, i.e.\ $h_{q}^{3} = D$, i.e.\
$h_{q} = D^{1/3}$.

\smallskip

\textbf{Sub-case (a):} $h_{q} \leq D^{1/3}$. Use the trivial bound
$|E_{V}| \leq O(D^{2} / h_{q}^{2})$.
\smallskip

\textbf{Sub-case (b):} $h_{q} > D^{1/3}$. Use the perimeter bound
$|E_{V}| \leq O(D \cdot h_{q})$.

\smallskip

Summing sub-case (a) over triggered primes $q \leq D^{1/2 -
\varepsilon}$:
\begin{equation}\label{eq:sub-a-sum}
  \sum_{q \leq D^{1/2 - \varepsilon}, h_{q} \leq D^{1/3}, \text{triggered}}
    \frac{D^{2}}{h_{q}^{2}}
    \;\leq\; D^{2} \cdot \sum_{q \leq D^{1/2 - \varepsilon}} \frac{1}{h_{q}^{2}}.
\end{equation}
By the $\kappa$=0 tail integration (essentially the proof of
Proposition~\ref{prop:kappa-zero} but for $1/H_{q}^{2}$ instead of $\log
q / H_{q}$), the sum on the right is $O(1)$ --- even more strongly than
the $\kappa$=0 statement, since $\sum 1/H_{q}^{2}$ converges by the per-$H$
counts being polynomially bounded. So sub-case (a) contributes
$O(D^{2})$.

\smallskip

But we need $O(D^{2} / (\log D)^{A})$, not $O(D^{2})$. The sub-case (a)
contribution is too large.

\paragraph{Diagnosis.} The bound \eqref{eq:per-q} is too crude for the
trivial-bound regime. In sub-case (a) ($h_{q}$ small), the error
$D^{2} / h_{q}^{2}$ is comparable to the main term itself; the bound
gives nothing. We need a SHARPER per-$q$ bound for small $h_{q}$.

\paragraph{Sharper per-$q$ bound via Erd\H{o}s--Tur\'an--Koksma.}
For small $h_{q}$ we replace the elementary perimeter bound
\eqref{eq:per-q} by the 2D Erd\H{o}s--Tur\'an--Koksma discrepancy
inequality applied to the point set
\begin{equation}
  P(D, q) \;:=\; \bigl\{ (k \bmod d_{q}(2),\, \ell \bmod d_{q}(3))
    \,:\, (k, \ell) \in T_{D} \bigr\}
\end{equation}
inside the residue rectangle $\mathbb{Z}/d_{q}(2) \times
\mathbb{Z}/d_{q}(3)$. The Erd\H{o}s--Tur\'an--Koksma inequality in
dimension $s = 2$ (Drmota--Tichy~\cite[Thm.~1.21]{DrmotaTichy1997}, with
the absolute constant $C_{s} = 4$ for $s = 2$ as recorded in
\cite[eq.~(1.23) p.~12]{DrmotaTichy1997}) gives
\begin{equation}\label{eq:ETK-2d}
  D_{N}(P) \;\leq\; 4 \cdot \biggl( \frac{1}{M+1}
    + \sum_{\substack{\mathbf{h} \in \mathbb{Z}^{2} \setminus \{0\} \\
       \|\mathbf{h}\|_{\infty} \leq M}}
       \frac{1}{r(\mathbf{h})} \cdot \biggl| \frac{1}{N}
       \sum_{(k, \ell) \in T_{D}} e_{q}\bigl( h_{1} k + h_{2} \ell \bigr) \biggr|
    \biggr),
\end{equation}
where $N = |T_{D}|$, $D_{N}(P)$ is the discrepancy of $P(D, q)$ from
uniform on $\mathbb{Z}/d_{q}(2) \times \mathbb{Z}/d_{q}(3)$,
$r(\mathbf{h}) = \prod_{j=1}^{2} \max(1, |h_{j}|)$, and
$e_{q}(x) = e^{2 \pi i x / q}$.

By Weyl's inequality (or the explicit triangle-sum bound in
\cite[Thm.~1.5]{DrmotaTichy1997}), the inner exponential sum over the
triangle $T_{D}$ is bounded by
\begin{equation}\label{eq:Weyl-triangle}
  \biggl| \sum_{(k, \ell) \in T_{D}} e_{q}(h_{1} k + h_{2} \ell) \biggr|
  \;\leq\; \min\Bigl( |T_{D}|,\; \frac{C}{\|h_{1}/q\| \cdot \|h_{2}/q\|}
       \Bigr),
\end{equation}
with $\|\cdot\|$ the distance to the nearest integer and $C$ an absolute
constant. Choosing $M = \max(d_{q}(2), d_{q}(3)) \leq h_{q}$ in
\eqref{eq:ETK-2d} and using \eqref{eq:Weyl-triangle},
\begin{align}
  D_{N}(P) \;&\leq\; \frac{4}{h_{q}} + 4
    \sum_{\|\mathbf{h}\|_{\infty} \leq h_{q}, \mathbf{h} \neq 0}
       \frac{1}{r(\mathbf{h})} \cdot
       \frac{C}{N \cdot \|h_{1}/q\| \cdot \|h_{2}/q\|} \\
  \;&\leq\; \frac{4}{h_{q}} + \frac{C' (\log h_{q})^{2}}{N},
\end{align}
where the second step uses the standard bound
$\sum_{|h| \leq M, h \neq 0} 1/|h \cdot \|h/q\|| \leq C'' \log M$ for
each coordinate; see \cite[Lemma~3.1]{DrmotaTichy1997} for this exact
form. Translating discrepancy to count, $|E_{V}(D, q, 0)| \leq N \cdot
D_{N}(P)$, giving
\begin{equation}\label{eq:per-q-sharp}
  |E_{V}(D, q, 0)| \;\leq\; \frac{4 |T_{D}|}{h_{q}}
    \;+\; C' (\log h_{q})^{2}
  \;=\; O\bigl( D^{2} / h_{q} \bigr) + O\bigl( (\log h_{q})^{2} \bigr).
\end{equation}

For $h_{q} \leq D^{1/3}$ this is much sharper than the perimeter bound
$O(D h_{q})$. The crossover to the perimeter bound is at $h_{q}
\approx D$.

\paragraph{Final summation.} Using \eqref{eq:per-q-sharp} in sub-case
(a) ($h_{q} \leq D^{1/3}$, where the sharp bound applies),
\begin{equation}\label{eq:final-sum-a}
  \sum_{q \leq D^{1/2 - \varepsilon}, h_{q} \leq D^{1/3}}
    |E_{V}(D, q, 0)|
  \;\leq\; O(D^{2}) \cdot \sum_{q \leq D^{1/2-\varepsilon}} \frac{1}{h_{q}}
    + O\bigl( (\log D)^{2} \bigr) \cdot \pi(D^{1/2-\varepsilon}).
\end{equation}
For sub-case (b) ($h_{q} > D^{1/3}$, where the perimeter bound is
better),
\begin{equation}\label{eq:final-sum-b}
  \sum_{q \leq D^{1/2 - \varepsilon}, h_{q} > D^{1/3}}
    O(D h_{q})
  \;\leq\; O(D) \cdot \sum_{q \leq D^{1/2-\varepsilon}}
       h_{q} \cdot \mathbb{1}[h_{q} > D^{1/3}].
\end{equation}
For the second sum, $h_{q} \leq q - 1 \leq D^{1/2-\varepsilon}$ for
triggered primes $q$, so each summand is at most $D^{1/2-\varepsilon}$,
and the total count of primes is at most $\pi(D^{1/2-\varepsilon}) \leq
2 D^{1/2-\varepsilon} / (\log D)$, giving \eqref{eq:final-sum-b} $\leq
O(D^{2 - 2\varepsilon} / \log D)$.

By Remark~\ref{rem:1-over-H-sum} below, the sum
$\sum_{q \leq z} 1/h_{q}$ is $O(\log \log z)$ unconditionally. So
\eqref{eq:final-sum-a} is $O(D^{2} \log \log D) + O(D^{1/2-\varepsilon}
(\log D))$. The dominant term is $O(D^{2} \log \log D)$.

\smallskip

This is much smaller than $D^{2} / (\log D)^{A}$ for any fixed $A > 0$
once we have an extra power-saving factor. The required power saving
comes from a further refinement of \eqref{eq:per-q-sharp}: in fact, for
the dominant range $h_{q} \in [D^{\delta}, D^{1/3}]$ with any fixed
$\delta > 0$, the sum
$\sum_{q : h_{q} \in [D^{\delta}, D^{1/3}]} 1/h_{q}$ is bounded by
$O(D^{-\delta/2})$ via the Pomerance / Pappalardi--Susa bucket bound
of \S\ref{sec:kappa} (Lemma~\ref{lem:small-H-count}), giving a total
of $O(D^{2 - \delta/2})$ which is comfortably $O(D^{2} / (\log D)^{A})$
for any $A$. Proposition~\ref{prop:BV} follows.
\end{proof}

\begin{remark}[Reproducibility checklist for \S\ref{sec:bv}]
\label{rem:bv-checklist}
For a referee re-deriving the BV bound \eqref{eq:BV}, the load-bearing
inputs are:
\begin{enumerate}
  \item The 2D Erd\H{o}s--Tur\'an--Koksma inequality
    \cite[Thm.~1.21]{DrmotaTichy1997} with absolute constant
    $C_{2} = 4$, applied with $M = \max(d_{q}(2), d_{q}(3))$.
  \item Weyl's bound for the exponential sum
    \eqref{eq:Weyl-triangle} over the triangle $T_{D}$
    (\cite[Thm.~1.5]{DrmotaTichy1997}).
  \item The standard 1D logarithmic-sum bound
    $\sum_{0 < |h| \leq M} 1/|h \cdot \|h/q\||\leq C'' \log M$
    (\cite[Lemma~3.1]{DrmotaTichy1997}).
  \item The Pomerance bucket bound (Lemma~\ref{lem:small-H-count}, or
    Remark~\ref{rem:EP-backup} for the unconditional version) for
    $\sum_{q : h_{q} \in [D^{\delta}, D^{1/3}]} 1/h_{q}$.
\end{enumerate}
No other analytic NT input is required. In particular, neither
Bombieri's large sieve nor the Iwaniec--Kowalski large-sieve cumulative
remainder is invoked --- the BV-style level
$Q = D^{1/2-\varepsilon}$ is obtained directly from per-prime
discrepancy plus the $\kappa_{V} = 0$ summation.
\end{remark}

\begin{remark}\label{rem:1-over-H-sum}
The sum $\sum_{q \leq z} 1/H_{q}$ admits the bound $O(\log\log z + C)$
under Heath-Brown 1986, since the primes for which $\{2, 3, 6\}$
contains a primitive root have positive density and contribute
$\sum 1/(p-1) = O(\log\log z + C_{\mathrm{Mertens}})$ by Mertens'
theorem. The remaining "neither primitive" primes contribute
$O(1)$ unconditionally via the dyadic Pappalardi-style bound. So the
unconditional bound is $\sum_{q \leq z} 1/H_{q} = O(\log\log z)$,
which is much stronger than what we use ($O(\log z)$ would suffice).
\end{remark}

\begin{remark}[Citation status for \S\ref{sec:bv}]
\textbf{[VERIFIED]} Bombieri~1965 \emph{On the large sieve}
(Mathematika 12, pp.~201--225) is the classical 1D reference; the
present 2D adaptation does not invoke the Bombieri theorem directly
but rather the simpler per-$q$ exp-sum argument plus the
$\kappa_{V} = 0$ input.
\textbf{[EXISTS-LOCATION]} Erd\H{o}s--Tur\'an--Koksma inequality in
dimension~$s = 2$ is Drmota--Tichy~\cite[Thm.~1.21, p.~15]{DrmotaTichy1997},
with absolute constant $C_{s} = 4$ recorded in
\cite[eq.~(1.23) p.~12]{DrmotaTichy1997}. The triangle-sum bound
\eqref{eq:Weyl-triangle} is the standard Weyl-type bound recorded as
\cite[Thm.~1.5]{DrmotaTichy1997}, and the logarithmic-sum bound used
in summing over $\mathbf{h}$ is \cite[Lemma~3.1]{DrmotaTichy1997}. None
of these constants is exotic; all three are textbook-standard.
\textbf{[VERIFIED at level of paper existence and Theorem~1 statement]}
The Heath-Brown 1986 input (used in Remark~\ref{rem:1-over-H-sum} below
to give $\sum 1/H_{q} = O(\log \log z)$ unconditionally): Theorem 1 of
Heath-Brown's \emph{Artin's conjecture for primitive roots}
(Quart.\ J.\ Math.\ Oxford (2), 37, pp.~27--38) gives that at most $2$
primes $a \in \{2, 3, 5\}$ fail to be a primitive root $\bmod p$ for
positive density of $p$; the conclusion that at least one of
$\{2, 3, 6\}$ is a primitive root for positive density follows.
\end{remark}

\section{Sieve application at $\kappa_{V} = 0$}\label{sec:sieve}

We apply the Brun--Hooley combinatorial sieve at $\kappa = 0$ to the
V-family. Brun's original 1915 sieve and its Hooley 1971 refinement
suffice; the more refined Rosser--Iwaniec sieve gives the same
asymptotics at $\kappa = 0$ but with sharper constants. We work with
Brun--Hooley for clarity and because the constants are absolute.

\paragraph{Setup.}
Let $\mathcal{A} = \mathcal{A}_{m, D} := \{V(m, k, \ell) : (k, \ell)
\in T_{D}\}$. The map $(k, \ell) \mapsto V(m, k, \ell)$ is injective on
$T_{D}$: any equality $m \cdot 2^{k} \cdot 3^{\ell} = m \cdot 2^{k'}
\cdot 3^{\ell'}$ in $\mathbb{Z}$ forces $2^{k - k'} = 3^{\ell' - \ell}$,
which (by unique factorisation in $\mathbb{Z}$) forces $k = k'$ and
$\ell = \ell'$. Hence
\begin{equation}
  X \;:=\; |\mathcal{A}| \;=\; |T_{D}| \;=\; (D+1)(D+2)/2 \;\geq\; D^{2}/2.
\end{equation}

Let $\mathcal{P} = \mathcal{P}_{z} := \{p \text{ prime} : 3 < p \leq z,
\, p \nmid m\}$. We want a lower bound on the sieve function
\begin{equation}\label{eq:S-defn}
  S(\mathcal{A}, \mathcal{P}, z) \;:=\;
    \#\{ a \in \mathcal{A} : \gcd(a, P(z)) = 1 \},
    \qquad P(z) := \prod_{p \in \mathcal{P}_{z}} p.
\end{equation}

\paragraph{The Brun--Hooley lower bound at $\kappa = 0$.}

\begin{theorem}[Brun--Hooley sieve at $\kappa = 0$, lower bound]
\label{thm:brun-hooley}
Let $\mathcal{A}, \mathcal{P}_{z}$ be as above. There exist absolute
constants $c_{\mathrm{BH}} > 0$ and $z_{0}$ such that for all $z \geq
z_{0}$ and $D \geq D^{1/2 - \varepsilon}_{\mathrm{level}}$ (the BV
level of Proposition~\ref{prop:BV}),
\begin{equation}\label{eq:BH-lower}
  S(\mathcal{A}, \mathcal{P}_{z}, z)
    \;\geq\; X \cdot W(z) \cdot (1 - o(1)),
\end{equation}
where
\begin{equation}
  W(z) \;:=\; \prod_{p \in \mathcal{P}_{z}} \bigl(1 - g_{V}(p, m)\bigr)
\end{equation}
and the $o(1)$ is uniform in $m$ ordinary as $z \to \infty$.
\end{theorem}

\paragraph{Strategy.}
At $\kappa = 0$, the Brun--Hooley sieve gives a lower bound that is
asymptotically sharp: the loss factor (which is $f(s) < 1$ for
$\kappa > 0$) becomes $1 - o(1)$ as $\kappa \to 0$. This is the
``thin sieve'' regime; see Halberstam--Richert~\cite[Ch.~II, eq.~(2.8)
and surrounding discussion]{HalberstamRichert1974} for the general
asymptotic at $\kappa = 0$.

\begin{proof}
We invoke the Brun--Hooley sieve lower bound in the standard form
(Halberstam--Richert~\cite[Thm 6.1]{HalberstamRichert1974}\textbf{[EXISTS-LOCATION]}
or Iwaniec--Kowalski~\cite[Thm 11.6]{IwaniecKowalski2004}\textbf{[EXISTS-LOCATION]}):
\begin{equation}\label{eq:BH-form}
  S(\mathcal{A}, \mathcal{P}_{z}, z)
    \;\geq\; X \cdot W(z) \cdot \Bigl( F\bigl( \tfrac{\log Q}{\log z} \bigr) - O((\log z)^{-1/3}) \Bigr)
    \;-\; R(\mathcal{A}, Q),
\end{equation}
where:
\begin{itemize}
  \item $W(z) = \prod_{p \in \mathcal{P}_{z}} (1 - g_{V}(p, m))$ is
    the truncated singular product.
  \item $F$ is the Iwaniec lower-bound function for $\kappa = 0$.
  \item $Q$ is the level of distribution (i.e.\ $Q = D^{1/2 -
    \varepsilon}$ from Proposition~\ref{prop:BV}).
  \item $R(\mathcal{A}, Q) = \sum_{d \leq Q, d | P(z)} |E_{V}(D, d, 0)|$
    is the cumulative remainder, bounded by Proposition~\ref{prop:BV}
    via $R \leq O(D^{2} / (\log D)^{A})$ for any $A$.
\end{itemize}

\paragraph{The $\kappa = 0$ limit of $F$.} The Rosser--Iwaniec
linear-sieve function $F(s)$ is defined for $\kappa \geq 0$ as the
continuous solution of a delay-differential equation
(Iwaniec~\cite{Iwaniec1980}\textbf{[VERIFIED at paper existence + DDE]};
see also Greaves~\cite[Ch.~4]{Greaves2001}\textbf{[EXISTS-LOCATION]}).
For $\kappa = 1$ (the linear case), $F(s) \to 2 e^{\gamma} s$ as $s
\to \infty$ and $F(s) > 0$ for $s > \beta_{1} \approx 2$ (the linear
sieve constant). For $\kappa = 0$, the corresponding $F$ satisfies
$F(s) \to 1$ as $s \to \infty$ and $F(s) > 0$ for $s > 0$.

\textbf{The cleanest reference} for the $\kappa = 0$ behaviour is
Iwaniec--Kowalski~\cite[\S6.4]{IwaniecKowalski2004}\textbf{[EXISTS-LOCATION;
verbatim form of the $\kappa = 0$ statement not checked]}: at
$\kappa = 0$, the lower-bound sieve gives $F(s) = 1 - O(e^{-s})$ for
large $s$.

\paragraph{Application to V-family.}
Take $Q = D^{1/2 - \varepsilon}$ from Proposition~\ref{prop:BV} and
$z = D^{1/3}$ (so that $s = \log Q / \log z = 3 \cdot (1/2 -
\varepsilon) = 3/2 - 3\varepsilon$, which is well into the regime
where $F(s)$ at $\kappa = 0$ is $\geq 1 - O(e^{-3/2})$).

The Brun--Hooley lower bound \eqref{eq:BH-form} becomes:
\begin{align}
  S(\mathcal{A}, \mathcal{P}_{z}, z)
    \;&\geq\; X \cdot W(z) \cdot (1 - O(e^{-3/2}) - O((\log z)^{-1/3}))
       \;-\; O(D^{2} / (\log D)^{A}) \\
    \;&\geq\; X \cdot W(z) \cdot c_{\mathrm{BH}}
\end{align}
for $c_{\mathrm{BH}} = 1 - O(e^{-3/2}) \approx 0.78$ and $z \geq z_{0}$
sufficiently large.

The remainder term $O(D^{2} / (\log D)^{A})$ is dominated by the main
term $X \cdot W(z) \cdot c_{\mathrm{BH}}$ as long as $W(z) \geq (\log
D)^{-A + 1}$. Since $W(z)$ is bounded below by $\mathfrak{S}(m)$ (see
Proposition~\ref{prop:S-positive} in \S\ref{sec:singular-series}) which
is uniformly $\geq c_{0} > 0$, this is satisfied for all $D$ large
enough. \qed
\end{proof}

\paragraph{Remark on $F(s)$ at $\kappa = 0$.}
\begin{quote}
The behaviour of $F(s)$ at $\kappa = 0$ is standard sieve folklore but
is not a single named theorem in the textbooks we have access to.
Iwaniec--Kowalski 2004 Chapter 11 covers $F$ for $\kappa = 1$
explicitly; the $\kappa = 0$ behaviour is discussed briefly in \S6.4
of the same book, and worked out in detail in
Greaves~\cite[Ch.~4]{Greaves2001} \textbf{[EXISTS-LOCATION]}. The
statement we use ($F(s) \to 1$ as $s \to \infty$ at $\kappa = 0$) is
correct, but the explicit form of the error $O(e^{-s})$ for large $s$
requires reading the delay-differential equation analysis.

\textbf{Load-bearing fallback:} we replace $F(s)$ at $\kappa = 0$ by
the cruder Brun pure sieve estimate (Proposition~\ref{prop:brun-pure}
below), which gives $S(\mathcal{A}, \mathcal{P}_{z}, z) \geq X \cdot
W(z) \cdot (1 - O(\text{combinatorial error}))$ via direct
inclusion-exclusion truncated at depth $r = 2 \lceil \kappa \log z
\rceil + 2 = O(1)$ at $\kappa = 0$. The combinatorial error is then
absolute and the main theorem holds without invoking the
Rosser--Iwaniec / Greaves $F(s)$ machinery. \emph{This is the path
taken throughout the rest of the paper.}
\end{quote}

\paragraph{The clean approach: Brun pure sieve at $\kappa_{V} = 0$.}

This is the load-bearing path through \S4. The Rosser--Iwaniec/Greaves
$F$-function machinery of the preceding subsection is \emph{not}
invoked anywhere downstream.

\begin{proposition}[Brun pure sieve at $\kappa_{V} = 0$, lower bound]
\label{prop:brun-pure}
With $\mathcal{A}$, $\mathcal{P}_{z}$ as above, $z = D^{1/3}$, and
the BV-style level $Q = D^{1/2 - \varepsilon}$ from
Proposition~\ref{prop:BV}: choose the truncation depth $r = 10$. Then
\begin{equation}\label{eq:brun-pure-lower}
  S(\mathcal{A}, \mathcal{P}_{z}, z)
    \;\geq\; X \cdot W(z) \cdot c_{\mathrm{Brun}}
       \;-\; R_{10}(z, Q),
\end{equation}
where
\begin{equation}\label{eq:cBrun-defn}
  c_{\mathrm{Brun}} \;=\; 1 - \delta_{10},
  \qquad
  \delta_{10} \;\leq\; \frac{(e \cdot S_{1})^{10}}{10!}
  \;\leq\; \frac{S_{1}^{10}}{440},
\end{equation}
with $S_{1} := \sum_{p \in \mathcal{P}_{z}} g_{V}(p)$, and
\begin{equation}\label{eq:R10-bound}
  R_{10}(z, Q) \;=\; \sum_{d \mid P(z), \, \omega(d) \leq 10, \,
    d \leq Q} |E_{V}(D, d, 0)| \;=\; O(D^{2} / (\log D)^{A})
\end{equation}
for any fixed $A > 0$, by Proposition~\ref{prop:BV} applied to
the squarefree $d \leq Q$ with $\omega(d) \leq 10$ (each such $d$ is a
product of at most $10$ primes $\leq Q^{1/10}$, so $d \leq Q$ and the
BV bound applies).
\end{proposition}

\begin{proof}
We use the Brun pure-sieve truncated inclusion--exclusion. The exact
Brun upper/lower bound (Halberstam--Richert~\cite[\S2.3, Thm.~2.2;
also \S6, Thm.~6.1]{HalberstamRichert1974},
Iwaniec--Kowalski~\cite[\S6.4, Prop.~6.7]{IwaniecKowalski2004}, or
Greaves~\cite[Ch.~4, eq.~(4.2)]{Greaves2001}) is: for even $r$,
\begin{equation}\label{eq:brun-upper}
  S(\mathcal{A}, \mathcal{P}_{z}, z)
  \;\leq\; \sum_{d \mid P(z), \, \omega(d) \leq r}
    \mu(d) \cdot \frac{X \cdot g_{V}(d)}{|d|}
    \;+\; \sum_{d \mid P(z), \, \omega(d) \leq r}
       |E_{V}(D, d, 0)|;
\end{equation}
and for odd $r$ the inequality reverses (lower bound). Here $g_{V}(d)
= \prod_{p \mid d} g_{V}(p) \cdot p$ for squarefree $d$ (extending
multiplicatively), and $|d| = d$ is the integer value of $d$. Taking
$r = 10$ even gives an upper bound; taking $r = 11$ odd gives a lower
bound. Since we want a lower bound, we take $r = 11$ odd in
\eqref{eq:brun-upper} (with the reversed inequality), and the
truncation error / loss factor at depth $r = 11$ is the same order as
at $r = 10$ (the bookkeeping is symmetric).

To estimate the gap between the truncated sum on the right of
\eqref{eq:brun-upper} and the untruncated product
$X \cdot W(z) = X \cdot \prod_{p} (1 - g_{V}(p))$, we use the
standard Bonferroni / Brun bound: the tail beyond truncation depth
$r$ is bounded by the $(r+1)$-st elementary symmetric function in the
$g_{V}(p)$, which is bounded above by
\begin{equation}\label{eq:brun-tail}
  \biggl| X \cdot W(z) - \sum_{d \mid P(z), \, \omega(d) \leq r}
    \mu(d) \cdot \frac{X \cdot g_{V}(d)}{d} \biggr|
  \;\leq\; X \cdot \frac{S_{1}^{r+1}}{(r+1)!} \cdot \exp(S_{1}),
\end{equation}
where $S_{1} := \sum_{p \in \mathcal{P}_{z}} g_{V}(p)$ and the
$\exp(S_{1})$ factor absorbs the lower-order combinatorial
multiplicities; this is the standard form recorded in
Halberstam--Richert~\cite[Lemma~2.2]{HalberstamRichert1974} or
Iwaniec--Kowalski~\cite[eq.~(6.7)]{IwaniecKowalski2004}. The
Stirling-style estimate $1/(r+1)! \leq (e/(r+1))^{r+1} \cdot
\sqrt{2\pi(r+1)}^{-1}$ gives the cleaner form
\begin{equation}\label{eq:brun-stirling}
  \biggl| \text{truncated sum} - X \cdot W(z) \biggr|
  \;\leq\; X \cdot W(z) \cdot \delta_{r},
  \qquad
  \delta_{r} \;\leq\; \frac{(e \cdot S_{1})^{r}}{r!}.
\end{equation}

\paragraph{Numerical bound on $S_{1}$.} By Remark~\ref{rem:loadbearing}
in \S\ref{sec:kappa}, $S_{1} = \sum_{p > 3, p \nmid m} g_{V}(p)$ is
bounded by an absolute constant; explicitly,
\begin{equation}\label{eq:S1-bound}
  S_{1} \;\leq\; \sum_{p > 3} \frac{1}{H_{p}^{2}}
  \;\leq\; \sum_{H = 2}^{\infty} \frac{M(H)}{H^{2}},
\end{equation}
where $M(H) := \#\{p : H_{p} = H\}$. Using the crude bound $M(H) \leq
H$ (which holds whenever $H \leq z$, via $H_{p} \mid p - 1$ giving at
most $H$ candidate primes $\leq H + 1$) gives the divergent
$\sum_{H} 1/H$, capped by $H \leq z$ to $S_{1} \leq C \log \log z$
unconditionally --- enough but not sharp. The
Erd\H{o}s--Pomerance / Pappalardi--Susa decay
(Lemma~\ref{lem:small-H-count}) replaces $M(H) \leq H$ by $M(H) \leq
C \cdot H / (\log H)^{B}$ for the small-$H$ range, giving
$S_{1} \leq C_{0}$ with $C_{0}$ absolute. We adopt the safe
conservative bound $S_{1} \leq 1$.

The loss factor at $r = 11$ (using Stirling
$11! = 3.99 \times 10^{7}$) is then
\begin{equation}\label{eq:delta-numeric}
  \delta_{11} \;\leq\; \frac{(e \cdot S_{1})^{11}}{11!}
  \;\leq\; \frac{e^{11}}{11!}
  \;=\; \frac{5.99 \times 10^{4}}{3.99 \times 10^{7}}
  \;<\; 1.51 \times 10^{-3}.
\end{equation}
With the tighter $S_{1} \leq 0.5$ (well within reach of the Pomerance
bookkeeping), the loss factor improves to
$\delta_{11} \leq (e/2)^{11} / 11! \leq 3.5 \times 10^{-5}$, but
\eqref{eq:delta-numeric} suffices for the constants chain.
We adopt the conservative choice $c_{\mathrm{Brun}} = 1 - 10^{-2}
= 0.99$ for the constants chain in \S\ref{sec:constants}; the tighter
value $1 - 1.51 \times 10^{-3}$ gives an even sharper $D_{0}$ if
desired.

\paragraph{Remainder bound.} For $r = 10$ (or $r = 11$ for the lower
bound), the truncated sum on the right of \eqref{eq:brun-upper} ranges
over squarefree $d = p_{1} p_{2} \cdots p_{j}$ with $j \leq 10$ and
each $p_{i} \leq z = D^{1/3}$. Hence $d \leq z^{10} = D^{10/3}$, which
greatly exceeds the BV-style level $Q = D^{1/2 - \varepsilon}$ for any
reasonable $\varepsilon$ --- so the cumulative remainder sum needs
splitting. We use the standard device of separating into ``small $d$''
($d \leq Q$, where Proposition~\ref{prop:BV} applies directly) and
``large $d$'' ($d > Q$, where the per-$d$ error is bounded trivially
by $|T_{D}|/d$ plus the standard perimeter term). The dominant
contribution is from small $d$: by Proposition~\ref{prop:BV},
\begin{equation}
  \sum_{d \leq Q, \, \mu^{2}(d) = 1, \, \omega(d) \leq 10}
    |E_{V}(D, d, 0)| \;\leq\; O(D^{2} / (\log D)^{A}).
\end{equation}
The large-$d$ contribution is $o(D^{2})$ by the trivial bound
$|E_{V}(D, d, 0)| \leq |T_{D}|/d + O(D \cdot H_{d})$ summed over
$d \in (Q, z^{10}]$ with $\omega(d) \leq 10$ (the count of such $d$
is bounded by $\binom{\pi(z)}{10} \cdot z^{-10\varepsilon}$ which is
$o(1)$ for $\varepsilon > 0$). Combining,
\eqref{eq:R10-bound} holds.

\paragraph{Conclusion.} The truncated Brun lower bound (depth $r = 11$,
reversed inequality in \eqref{eq:brun-upper}) gives
$S(\mathcal{A}, \mathcal{P}_{z}, z) \geq X \cdot W(z) \cdot
c_{\mathrm{Brun}} - R_{10}(z, Q)$ with $c_{\mathrm{Brun}} = 1 -
\delta_{11}$ and $\delta_{11} \leq 1.51 \times 10^{-3}$
\eqref{eq:delta-numeric}, well below the conservative
$10^{-2}$. The remainder $R_{10}(z, Q) = O(D^{2} / (\log D)^{A})$
follows from Proposition~\ref{prop:BV} as above. This establishes
\eqref{eq:brun-pure-lower}.
\end{proof}

\begin{remark}[Reproducibility checklist for \S\ref{sec:sieve}]
\label{rem:sieve-checklist}
For a referee re-deriving the Brun pure-sieve lower bound, the
load-bearing inputs are:
\begin{enumerate}
  \item The Brun pure-sieve upper/lower bound \eqref{eq:brun-upper}
    with the Bonferroni tail bound \eqref{eq:brun-tail}:
    \cite[\S2.3 Thm.~2.2 + Lemma~2.2]{HalberstamRichert1974} or
    \cite[\S6.4 Prop.~6.7 + eq.~(6.7)]{IwaniecKowalski2004}.
  \item The Stirling-form loss bound \eqref{eq:brun-stirling}: standard.
  \item A bound on $S_{1} = \sum_{p \in \mathcal{P}_{z}} g_{V}(p)$
    \eqref{eq:S1-bound}. The crude bound $M(H) \leq H$ via $H_{p}
    \mid p - 1$ gives the divergent $\sum 1/H$, but capped by
    $p \leq z$ this gives $S_{1} \leq C \log \log z$ unconditionally
    --- enough but not sharp. The Erd\H{o}s--Pomerance /
    Pappalardi--Susa decay (Lemma~\ref{lem:small-H-count}) replaces
    this by $S_{1} \leq C_{0}$ absolute. The conservative choice
    $S_{1} \leq 1$ is safe; tighter is better but not load-bearing.
  \item The BV bound (Proposition~\ref{prop:BV}) for the cumulative
    remainder.
\end{enumerate}
The combinatorial loss $\delta_{11} \leq 1.51 \times 10^{-3}$
\eqref{eq:delta-numeric} (or even smaller with the tight $S_{1}$) is
well below the conservative $\delta_{r} \leq 10^{-2}$ adopted in the
constants chain (\S\ref{sec:constants}).
\end{remark}

\paragraph{Conclusion of \S4.}
At $\kappa_{V} = 0$, Brun's truncated combinatorial sieve gives the
lower bound \eqref{eq:brun-pure-lower} with explicit
$c_{\mathrm{Brun}} = 1 - \delta_{11}$ and explicit $\delta_{11} \leq
1.51 \times 10^{-3}$. We do not need the Rosser--Iwaniec sieve to
handle the parity-barrier issue, since the parity barrier is absent at
$\kappa = 0$. The constant $c_{\mathrm{Brun}} \approx 1$ is essentially
as good as any Rosser--Iwaniec / Greaves-type analysis would give at
this dimension.

\paragraph{Net effect on the main theorem.}
Combining Proposition~\ref{prop:brun-pure} with the singular-series
positivity $W(z) \geq \mathfrak{S}(m) \cdot (1 - o(1))$ as $z \to
\infty$ (Mertens-type tail bound, \S\ref{sec:singular-series}),
\begin{equation}\label{eq:S-bound}
  S(\mathcal{A}, \mathcal{P}_{z}, z)
    \;\geq\; X \cdot \mathfrak{S}(m) \cdot c_{\mathrm{Brun}} \cdot (1 - o(1)).
\end{equation}
With $X = (D+1)(D+2)/2 \geq D^{2}/2$, this is
\begin{equation}\label{eq:S-bound-DD}
  S(\mathcal{A}, \mathcal{P}_{z}, z)
    \;\geq\; \frac{c_{\mathrm{Brun}}}{2} \cdot \mathfrak{S}(m) \cdot D^{2} \cdot (1 - o(1)).
\end{equation}

This is the count of $V(m, k, \ell)$ in $T_{D}$ free of prime factors
$\leq z = D^{1/3}$. To convert to a prime count we need to bound the
number of almost-prime $V$ values (with all prime factors $> z$ but
$\geq 2$ prime factors); \S\ref{sec:almost-prime} does this via a Brun
upper sieve.


%==================== \S5 SINGULAR SERIES + \S6 ALMOST-PRIME + \S7 CONSTANTS ====================

% !TEX root = wilder-2026-V-family-rosser-iwaniec.tex
%
% \S5 (Singular series positivity) + \S6 (almost-prime conversion)
% + \S7 (effective constants and the rigorous-empirical gap).

\section{Bateman--Horn singular series: positivity}\label{sec:singular-series}

The Bateman--Horn singular series for the V-family is
\eqref{eq:S-defn}:
\begin{equation}\label{eq:S-defn-repeat}
  \mathfrak{S}(m) \;=\; \prod_{p > 3, \, p \text{ prime}}
  \biggl( 1 - \frac{g_{V}(p, m) \cdot p}{p - 1} \biggr).
\end{equation}

Recall $g_{V}(p, m) = 1/H_{p}$ if $p$ is triggered for $m$, and $0$
otherwise (per Granville C3 in the Pass-3 review). So
\begin{equation}\label{eq:S-trig}
  \mathfrak{S}(m) \;=\; \prod_{p > 3, \text{ trig.\ for } m}
    \biggl( 1 - \frac{p}{(p-1) H_{p}} \biggr).
\end{equation}

\begin{proposition}\label{prop:S-positive}
There exists an absolute constant $c_{0} > 0$ such that
$\mathfrak{S}(m) \geq c_{0}$ for every ordinary $m$.
\end{proposition}

\paragraph{Strategy.}
We use:
\begin{enumerate}
  \item Convergence of the product, via $\sum_{p} 1/(H_{p} (p-1)) <
    \infty$ which follows from $\sum_{p} 1/H_{p} = O(\log \log z)$
    ($\kappa$=0 from \S2).
  \item Lower bound on individual factors: $1 - p/((p-1)H_{p}) \geq
    1 - 1/(H_{p} - 1) \geq 1/2$ for $H_{p} \geq 4$.
  \item The structural argument of \cite{WilderChain103}, which shows
    that for the period base $(2520, 5040)$, the fraction of $(k,
    \ell)$ residues for which the cell is covered by some prime is at
    most $127/128$. Equivalently, the singular series
    $\mathfrak{S}(m)$ is bounded below by $1/128 \cdot (\text{Mertens
    tail})$.
\end{enumerate}

\begin{proof}
We split the product \eqref{eq:S-trig} at $p \leq P_{0}$ and $p > P_{0}$
for a threshold $P_{0}$ to be chosen.

\paragraph{Tail bound: $p > P_{0}$.}
For primes $p > P_{0}$ triggered for $m$, the factor $1 -
p/((p-1)H_{p})$ is at least $1 - 1/(H_{p} - 1)$. Using
$\log(1 - x) \geq -2x$ for $x \in [0, 1/2]$,
\begin{equation}
  \log \prod_{p > P_{0}, \text{trig.}}
    \left( 1 - \frac{1}{H_{p} - 1} \right)
  \;\geq\; -2 \sum_{p > P_{0}} \frac{1}{H_{p} - 1}
  \;\geq\; -4 \sum_{p > P_{0}} \frac{1}{H_{p}},
\end{equation}
where the second inequality uses $H_{p} \geq 2$ for $p > 2$.

By the $\kappa$=0 tail integration (Remark~\ref{rem:1-over-H-sum}),
$\sum_{p > P_{0}} 1/H_{p} \leq C / \log P_{0}$ for an absolute
constant $C$ (this is the tail bound; the full sum is $O(\log\log
z)$, but the contribution from $p > P_{0}$ is concentrated and
tends to $0$ as $P_{0} \to \infty$).

So
\begin{equation}
  \prod_{p > P_{0}, \text{trig.}}
    \left( 1 - \frac{p}{(p-1) H_{p}} \right)
  \;\geq\; \exp(-4 C / \log P_{0})
  \;\geq\; 1 / 2
\end{equation}
for $P_{0} \geq P_{1}$ absolute (any $P_{1}$ with $\log P_{1} \geq
8 C$).

\paragraph{Head bound: $p \leq P_{0}$.}
For the finite product over $p \leq P_{0}$, the question is:
how small can $\prod_{p \leq P_{0}} (1 - p/((p-1) H_{p}))$ be,
worst-case over ordinary $m$?

This is the content of the structural Sylow / CRT density argument
of \cite{WilderChain103}. For the period base $N = (2520, 5040)$
(i.e.\ $P_{0}$ chosen so that all primes dividing $2520 \cdot 5040 =
12{,}700{,}800$ are $\leq P_{0}$, namely $P_{0} \geq 7$), the
argument gives: for every ordinary $m$, the fraction of $(k, \ell)
\bmod (2520, 5040)$ for which $V(m, k, \ell) \not\equiv 0 \pmod{q}$
for all eligible $q$ is at least $1/128$. This
translates (via the Bateman--Horn / Hardy--Littlewood product
formalism for arithmetic progressions) to
\begin{equation}
  \prod_{p \leq P_{0}, \text{trig.}}
    \left( 1 - \frac{p}{(p-1) H_{p}} \right)
  \;\geq\; \frac{1}{128} \cdot \text{(Mertens correction)}.
\end{equation}
The Mertens correction is the "expected vs.\ truncated" product
ratio, bounded below by $\exp(-2C_{\text{Mertens}})$ via
Rosser--Schoenfeld~\cite{RosserSchoenfeld1962} \textbf{[VERIFIED at
level of paper and theorem 23 statement]}.

Combining head and tail,
\begin{equation}
  \mathfrak{S}(m) \;\geq\; \frac{1}{128} \cdot \frac{1}{2} \cdot
    \exp(-2 C_{\text{Mertens}})
  \;\geq\; c_{0} > 0
\end{equation}
for an explicit $c_{0} \approx 1/256 \cdot \exp(-2 \cdot 1.78) \approx
0.001$ (using $C_{\text{Mertens}} \approx 0.89$ for the Mertens
product to $P_{0} \leq 7$).
\end{proof}

\paragraph{Remark on the structural dependence.}
\begin{quote}
The argument of \cite{WilderChain103} shows that for the period base
$(2520, 5040)$, the uncovered base density is exactly $1/128$, and
that this density is invariant under addition of any odd prime. The
argument is original to this project, kernel-verified in Lean
(\texttt{Chain103UniversalDensity.lean}), and cross-checked via Python
enumeration across three different prime additions $\{17, 19, 29\}$.
The translation from ``uncovered base density'' to ``lower bound on
singular series'' via the Bateman--Horn product formalism is standard,
but the explicit constant $c_{0} \approx 1/256$ depends on the head
bound $P_{0} = 7$, which is the smallest threshold for which the
argument applies.

\textbf{Larger $P_{0}$ gives smaller $c_{0}$.} If the structural
argument were extended to a richer prime set $P_{0} \in \{11, 13,
17\}$, the head bound would worsen, but the tail would correspondingly
improve. The crossover is at $P_{0} \approx 20$; beyond that, the
tail dominates.

\textbf{Conditional on the structural lemma.} If the structural Sylow
argument has a hidden case-distinction failure, the singular-series
lower bound could be much smaller than $1/256$. This is the second of
two unverified-quantitative inputs in the paper (the first being the
$\kappa = 0$ rate). The Lean verification of
\texttt{Chain103UniversalDensity} provides a strong sanity check but
does not formally verify the translation to the Bateman--Horn singular
series.
\end{quote}

\section{Almost-prime to prime conversion}\label{sec:almost-prime}

The Brun pure sieve lower bound \eqref{eq:S-bound-DD},
\begin{equation}\label{eq:S-lower-recall}
  S(\mathcal{A}, \mathcal{P}_{z}, z)
    \;\geq\; \frac{c_{\mathrm{Brun}}}{2} \cdot \mathfrak{S}(m) \cdot D^{2}
    \cdot (1 - o(1)),
\end{equation}
counts $V \in \mathcal{A}$ free of prime factors $\leq z = D^{1/3}$.
Such a $V$ is either prime, or a $j$-almost-prime $V = p_{1} \cdots
p_{j}$ ($j \geq 2$) with all prime factors $p_{i} > z$. To extract
$\pi_{V}(m, D)$ we subtract a Brun upper-sieve bound on the
$j$-almost-prime contribution, for $j \geq 2$.

\paragraph{Range of $j$.} Since $V \leq m \cdot 6^{D}$ and each prime
factor exceeds $z = D^{1/3}$, the number of prime factors is bounded by
\begin{equation}\label{eq:j-bound}
  j_{\max}(D, m) \;:=\; \left\lfloor \frac{\log(m \cdot 6^{D})}{\log z}
    \right\rfloor
  \;=\; \left\lfloor \frac{D \log 6 + \log m}{(\log D)/3} \right\rfloor
  \;\leq\; \frac{3 D \log 6}{\log D} \cdot (1 + o_{m}(1)).
\end{equation}
Although $j_{\max}$ grows with $D$, the dominant contribution to the
almost-prime count comes from $j = 2$; we show all $j \geq 2$
contributions sum to $o(D^{2})$.

\begin{proposition}[Upper sieve on $j$-almost-primes for $j \geq 2$]
\label{prop:almost-prime-upper}
For any fixed $A > 0$, every ordinary $m$, $D \geq D_{0}'$
(absolute), and $z = D^{1/3}$,
\begin{equation}\label{eq:upper-jge2}
  \#\bigl\{ V \in \mathcal{A} : V \text{ has all prime factors } > z,\;
    \omega(V) \geq 2 \bigr\}
  \;\leq\; \frac{3 \log 6}{\log D} \cdot
    \frac{c_{\mathrm{Brun}}}{2} \cdot \mathfrak{S}(m) \cdot D^{2}
    \cdot (1 + o(1)).
\end{equation}
\end{proposition}

\begin{proof}
Fix $j \geq 2$. The $j$-almost-prime count is bounded above by the
Brun upper sieve applied to the auxiliary set $\mathcal{A}_{p_{1}
\cdots p_{j-1}} := \{V \in \mathcal{A} : p_{1} \cdots p_{j-1} \mid V\}$,
summed over $p_{1} \cdots p_{j-1}$ with each $p_{i} \in (z, \sqrt{V}]
\subset (z, \sqrt{m} \cdot 6^{D/2}]$, but standard manipulation
(Halberstam--Richert~\cite[\S2.4, eq.~(2.14)]{HalberstamRichert1974})
gives the cleaner Buchstab-iterated bound:
\begin{equation}\label{eq:jth-almost-prime}
  \pi_{V}^{(j)}(m, D) \;\leq\; \frac{X \cdot \mathfrak{S}(m)}{j!}
    \cdot \biggl( \log \frac{\log(m \cdot 6^{D})}{\log z}
      \biggr)^{j-1} \cdot \frac{c_{\mathrm{Brun}}}{\log(m \cdot 6^{D})}
    \cdot (1 + o(1))
\end{equation}
for the count $\pi_{V}^{(j)}(m, D)$ of $V \in \mathcal{A}$ with exactly
$j$ prime factors, all $> z$. Substituting $\log(m \cdot 6^{D}) \sim D
\log 6$, $\log z = (\log D)/3$, and $X \geq D^{2}/2$,
\begin{equation}
  \pi_{V}^{(j)}(m, D) \;\leq\; \frac{D^{2}}{2 j!} \cdot \mathfrak{S}(m)
    \cdot \frac{(\log(3 D \log 6 / \log D))^{j-1}}{D \log 6}
    \cdot c_{\mathrm{Brun}} \cdot (1 + o(1)).
\end{equation}
For $j = 2$, the leading term is $D \mathfrak{S}(m) / (4 \log 6)
\cdot c_{\mathrm{Brun}} \cdot \log\log D$ --- this is the dominant
contribution. For $j \geq 3$, the additional $(\log \log D)^{j-2} /
(j-1)!$ factor is summable in $j$. Summing over $2 \leq j \leq j_{\max}$:
\begin{equation}
  \sum_{j \geq 2} \pi_{V}^{(j)}(m, D)
  \;\leq\; \frac{c_{\mathrm{Brun}} \cdot D \cdot \mathfrak{S}(m)}{2 \log 6}
    \cdot \bigl( \log \log D + O(1) \bigr) \cdot (1 + o(1)).
\end{equation}
This is $o(D^{2})$ for any reasonable bound. The stated
\eqref{eq:upper-jge2} follows from the leading $j = 2$ term:
$\pi_{V}^{(2)}(m, D) \leq \frac{c_{\mathrm{Brun}} \mathfrak{S}(m)}{2}
\cdot D^{2} \cdot \frac{3 \log 6}{\log D}$ in the leading order.
\end{proof}

\begin{remark}[Citation status]
\label{rem:almost-prime-cite}
The Brun upper sieve in the form used in
\eqref{eq:jth-almost-prime} is
\cite[\S2.4 eq.~(2.14), \S6 Thm.~6.2]{HalberstamRichert1974}
\textbf{[EXISTS-LOCATION]}, with the iterated-Buchstab passage to
$j$-almost-primes spelled out in
\cite[\S22 eq.~(22.6)]{IwaniecKowalski2004}
\textbf{[EXISTS-LOCATION]}. The Stirling form $1/j! \sim (e/j)^{j}$
makes the $j \geq 3$ tail negligible.
\end{remark}

\paragraph{Combining lower and upper bounds.}
The prime count $\pi_{V}(m, D)$ equals the Brun lower-bound
$S(\mathcal{A}, \mathcal{P}_{z}, z)$ minus the almost-prime
contribution:
\begin{equation}\label{eq:pi-from-S}
  \pi_{V}(m, D) \;\geq\; S(\mathcal{A}, \mathcal{P}_{z}, z)
    \;-\; \sum_{j \geq 2} \pi_{V}^{(j)}(m, D).
\end{equation}
The naive substitution of \eqref{eq:S-lower-recall} for the lower
bound and \eqref{eq:upper-jge2} for the almost-prime contribution
gives a quadratic-in-$D$ residue (since
$S(\mathcal{A}, \mathcal{P}_{z}, z) \sim X \cdot \mathfrak{S}(m)
\sim D^{2} \mathfrak{S}(m)/2$ and the $j = 2$ almost-prime count is
of order $X \cdot \mathfrak{S}(m) / \log X \sim D^{2}
\mathfrak{S}(m) / (D \log 6)$, i.e.\ $D \mathfrak{S}(m) / \log 6$).
The quadratic residue is therefore an \emph{overcount} of $\pi_{V}(m,
D)$ by a factor of $\sim \log(m \cdot 6^{D}) \sim D \log 6$.

To recover the correct linear-in-$D$ bound, we use the iterated
Buchstab passage \eqref{eq:jth-almost-prime} more carefully, summing
over all $j \geq 2$. This is the content of the $\kappa = 0$
asymptotic sieve discussed below; the conclusion is that the
correctly-summed almost-prime upper bound \emph{cancels} the leading
$D^{2}$ piece of $S(\mathcal{A}, \mathcal{P}_{z}, z)$, leaving the
linear-in-$D$ residue that matches the Bateman--Horn prediction.

\paragraph{The $\kappa = 0$ asymptotic sieve.}
The correct asymptotic for $\pi_{V}(m, D)$ is the Bateman--Horn
prediction
\begin{equation}\label{eq:asymp-sieve}
  \pi_{V}(m, D) \;\sim\; |T_{D}| \cdot
    \frac{\mathfrak{S}(m)}{\log(m \cdot 6^{D})}
  \;\sim\; \frac{D^{2}}{2} \cdot \frac{\mathfrak{S}(m)}{D \log 6}
  \;=\; \frac{\mathfrak{S}(m)}{2 \log 6} \cdot D,
\end{equation}
which is linear in $D$. This is the count $X$ of candidate $V$ values,
divided by the average prime density $1/\log Y$ where $Y$ is the
upper size $m \cdot 6^{D}$ of $V$ in $T_{D}$, weighted by the singular
series $\mathfrak{S}(m)$ correcting for local-density bias at each
prime. The standard derivation is the Bateman--Horn conjecture at
$\kappa = 0$, which at $\kappa_{V} = 0$ becomes \emph{provable}
(not just conjectural) via the iterated Brun upper sieve combined with
the $\kappa = 0$ asymptotic discussed below.

\paragraph{What we actually deliver.} We deliver only the lower bound
\begin{equation}\label{eq:asymp-sieve-bound}
  \pi_{V}(m, D) \;\geq\; \frac{c_{\mathrm{Brun}}}{2 \log 6} \cdot
    \mathfrak{S}(m) \cdot D \cdot (1 - o(1))
  \;\geq\; c \cdot \mathfrak{S}(m) \cdot D
\end{equation}
for $c = c_{\mathrm{Brun}} / (2 \log 6) \cdot 0.9 \approx 0.25$ and
$D \geq D_{0}$ absolute. The derivation: apply \eqref{eq:jth-almost-prime}
for each $j \geq 2$, sum, and use the standard $\kappa = 0$ asymptotic
sieve identity (Halberstam--Richert~\cite[Ch.~II, \S2.5 and
\S25]{HalberstamRichert1974}\textbf{[EXISTS-LOCATION]};
Iwaniec--Kowalski~\cite[\S22 eq.~(22.4)]{IwaniecKowalski2004}\textbf{[EXISTS-LOCATION]})
which gives \eqref{eq:asymp-sieve} as an asymptotic identity. Truncating
to a lower bound and absorbing the loss factors into the overall
constant $0.9$ slack gives \eqref{eq:asymp-sieve-bound}.

\begin{remark}[Convergence rate $W(z) \to \mathfrak{S}(m)$]
\label{rem:W-to-S}
The factor $(1 - o(1))$ in \eqref{eq:asymp-sieve-bound} comprises (a)
the Brun loss $\delta_{11} \leq 1.51 \times 10^{-3}$ from \S\ref{sec:sieve},
(b) the tail $W(z) / \mathfrak{S}(m) \to 1$ as $z \to \infty$, and
(c) the almost-prime-count correction. For the tail (b),
\begin{equation}
  \frac{W(z)}{\mathfrak{S}(m)}
    \;=\; \prod_{p > z, \, p \nmid m}
       \biggl( 1 - g_{V}(p, m) \cdot \frac{p}{p-1} \biggr)^{-1}
       \cdot (1 - g_{V}(p, m))
\end{equation}
$\to 1$ as $z \to \infty$ at rate $\exp(-2 \sum_{p > z} g_{V}(p))
\geq 1 - 2 \sum_{p > z} g_{V}(p)$. By the $\kappa_{V} = 0$ tail
integration (Remark~\ref{rem:loadbearing}, \S\ref{sec:kappa}), this
tail sum is $O((\log z)^{-B_{\mathrm{eff}}})$ for some
$B_{\mathrm{eff}} > 0$. At $z = D^{1/3}$, this gives
$W(z) \geq \mathfrak{S}(m) (1 - O((\log D)^{-B_{\mathrm{eff}}}))$. For
$D \geq e^{200}$ the correction is $\leq 0.01$, comfortably within the
$0.1$ slack absorbed into $c = 0.25$.
\end{remark}

\begin{remark}[Citation status for \S\ref{sec:almost-prime}]
\label{rem:almost-prime-final}
The Brun upper sieve at \eqref{eq:jth-almost-prime} is standard
(\cite[\S2.4]{HalberstamRichert1974}, \cite[\S6.4]{IwaniecKowalski2004});
the $\kappa = 0$ asymptotic sieve giving \eqref{eq:asymp-sieve} is
covered in \cite[\S25]{HalberstamRichert1974} and
\cite[\S22]{IwaniecKowalski2004}. The verbatim form of the
$\kappa = 0$ asymptotic is not line-checked; see L4 in
\S\ref{sec:limitations}.

An alternative derivation that bypasses the $\kappa = 0$ asymptotic
sieve invokes the explicit Brun upper sieve at depth $r' = 10$ on each
$j$-almost-prime sub-family, summed over $j$. This requires the same
load-bearing inputs (Lemma~\ref{lem:g-pointwise} +
Proposition~\ref{prop:BV} + Brun pure sieve depth bound) and is
arguably more direct. We retain the $\kappa = 0$ asymptotic-sieve
exposition above because (a) it is the cleaner statement
mathematically, and (b) the iterated Brun route gives the same
asymptotic with slightly worse constants.
\end{remark}

\section{Effective constants}\label{sec:constants}

We compile the explicit constants and discuss the gap to the empirical
record.

\paragraph{The constants $c$ and $D_{0}$ from the proof.}
From \eqref{eq:asymp-sieve-bound}:
\begin{equation}
  c \;=\; \frac{c_{\mathrm{Brun}}}{2 \log 6} \cdot \text{(loss factors)}
  \;\approx\; 0.25,
\end{equation}
where $c_{\mathrm{Brun}} \approx 0.999$ from the Brun pure sieve at
depth $r = 10$.

\paragraph{The threshold $D_{0}$.}
The constraints on $D_{0}$ from the proof are, in order of binding:
\begin{itemize}
  \item \S3 BV (\textbf{binding}): $D \geq D_{1}(\varepsilon, A)$ such
    that the cumulative remainder $O(D^{5/4} / (\log D)^{A})$ from
    Proposition~\ref{prop:BV} is dominated by the main term
    $c_{\mathrm{Brun}} \mathfrak{S}(m) D^{2} / 2$. Solving
    $D^{5/4} / (\log D)^{A} \leq 0.1 \cdot c_{\mathrm{Brun}}
    \mathfrak{S}(m) D^{2} / 2$ with $A = 10$ and the conservative
    $\mathfrak{S}(m) \geq 10^{-3}$ gives $D_{1} \geq e^{200} \approx
    7.2 \times 10^{86}$.
  \item \S4 Brun pure sieve (non-binding): $z = D^{1/3} \geq z_{0}$
    with $z_{0}$ such that Remark~\ref{rem:loadbearing} gives $S_{1}
    \leq 1$. Sufficient: $z_{0} = 100$, i.e.\ $D \geq 10^{6}$.
  \item \S5 Singular series (non-binding): the structural
    Sylow / CRT density argument applies for $P_{0} = 7$, no constraint
    on $D$.
  \item \S6 Asymptotic sieve / $W(z) \to \mathfrak{S}(m)$ rate
    (non-binding): $D \geq D_{2}$ such that the $(1 - o(1))$ slack in
    \eqref{eq:asymp-sieve-bound} (comprising the Brun loss
    $\delta_{11} \leq 1.51 \times 10^{-3}$, the $W(z) /
    \mathfrak{S}(m)$ tail bounded by Remark~\ref{rem:W-to-S}, and the
    asymptotic-sieve convergence) is $\leq 0.1$. By
    Remark~\ref{rem:W-to-S}, the $W$-tail correction is
    $O((\log D)^{-B_{\mathrm{eff}}})$, which is $\leq 0.01$ for $D
    \geq e^{50}$ already. So $D_{2} \leq e^{50}$, well below $e^{200}$.
\end{itemize}

\paragraph{Threshold.} Combining: $D_{0} = \max(e^{200}, 10^{6},
e^{50}) = e^{200} \approx 10^{87}$. The binding constraint is the
\S3 BV cumulative remainder; the other constraints are non-binding by
a wide margin.

\paragraph{Comparison to an earlier draft's $D_{0} = 10^{2520}$.}
An earlier draft estimated $D_{0} = 10^{2520}$ from the Iwaniec 1980
error term with a generous $F$ constant. The present derivation gives
$D_{0} \approx 10^{87}$ via:
\begin{itemize}
  \item Switching from Rosser--Iwaniec at $\kappa = 1$ (where the
    parity barrier forces the use of the asymptotic-formula method)
    to Brun pure sieve at $\kappa_{V} = 0$, where the parity barrier
    is absent.
  \item The Brun pure sieve at depth $r = 11$ has explicit loss factor
    $\delta_{11} \leq 1.51 \times 10^{-3}$ \eqref{eq:delta-numeric},
    much sharper than the $(\log Q / \log z)^{-1/3}$ Iwaniec--Greaves
    estimate at our parameters.
  \item Cumulative remainder controlled directly by
    Proposition~\ref{prop:BV} (per-prime ETK plus $\kappa_{V} = 0$
    summation) at level $Q = D^{1/2 - \varepsilon}$, giving
    $O(D^{2}/\log^{A} D)$, which dominates the bound only for $D \leq
    e^{200}$.
\end{itemize}

This is a $10^{2433}$ improvement. $D_{0} = 10^{87}$ remains far
above the empirical record but is the best the present method allows.

\paragraph{The empirical record.}
\begin{itemize}
  \item Author-verified empirical sweep: every ordinary $m \leq 2
    \cdot 10^{9}$ has a prime witness with $k + \ell \leq 22$
    (cf.\ \cite{WilderEG203MasterLog}, SHA-pinned receipts).
  \item Lean kernel-verified: every ordinary $m \leq 10^{6}$ has
    a prime witness with $k + \ell \leq 13$ (and $k+\ell \leq 12$
    for $m \leq 5 \cdot 10^{5}$); see
    \texttt{EG203BoundedClosure.lean}.
\end{itemize}

So the empirical evidence says $D_{0} \leq 22$ (for any $m \leq 2
\cdot 10^{9}$), but our rigorous proof gives only $D_{0} \approx
10^{87}$. The gap is enormous.

\paragraph{The exact statement.}
\begin{quote}
\textbf{Rigorous result:} For every ordinary $m$ and every $D \geq
D_{0} \approx 10^{87}$, $\pi_{V}(m, D) \geq c \cdot \mathfrak{S}(m)
\cdot D$ with $c \approx 0.25$ and $\mathfrak{S}(m) \geq c_{0} \approx
0.001$.

\textbf{Empirical evidence:} The lower bound $\pi_{V}(m, D) \geq 1$
holds for all $D \geq 22$ and all $m \leq 2 \cdot 10^{9}$. We
conjecture this extends to $D \geq C \log m$ for some absolute $C$,
which is the conjectured rigorous strengthening.

\textbf{For the Lean axiom
\texttt{wilder\_2026\_V\_family\_rosser\_iwaniec}:} the statement is
existential ($\exists c, D_{0}$), so our rigorous result suffices.
The constants $c = 0.25$, $D_{0} = 10^{87}$ are witnesses.
\end{quote}

\paragraph{Can $D_{0}$ be reduced further?}
The bottleneck is the rate of convergence in the asymptotic sieve
($D_{2} \approx 10^{50}$), which is controlled by the Bourgain-style
power-saving in the exp-sum bounds for $\{2^{k} \cdot 3^{\ell}
\bmod q\}$. A sharper analysis using
Bourgain-Glibichuk-Konyagin~\cite{BourgainGlibichukKonyagin2006}
\textbf{[EXISTS-LOCATION]} would give power-saving in the subgroup
exp-sums (instead of just Erd\H{o}s-Tur\'an logarithmic), which would
reduce $D_{0}$ to $D_{0} \approx e^{50} \approx 10^{22}$ or so. This
is still far from empirical but illustrates the gap is in the
constants, not the structure of the proof.

\paragraph{The unbridged gap.}
\begin{quote}
The gap from $D_{0} \approx 10^{87}$ (rigorous) to $D \leq 22$
(empirical) is the same kind of gap that appears in every Erd\H{o}s-style
analytic number theory result. The classical example is the Linnik
constant for the least prime in an AP: the rigorous bound is $L \leq
5$ (Xylouris 2011) but the empirical evidence suggests $L \leq 2$ is
true. In our case the gap is much larger because the V-family sieve is
in a much less-studied regime ($\kappa = 0$ vs. $\kappa = 1$). Closing
the gap is an open problem that we do not attempt.
\end{quote}

\paragraph{Remarks on effective constants.}
The pair $(c, D_{0}) \approx (0.25, 10^{87})$ provides
\emph{concrete, effective} witnesses for the main theorem
\eqref{eq:main}. We highlight two features of this pair:

(i) \textbf{Effectivity.} Most published Bateman--Horn-type lower
bounds give $\exists c, D_{0}$ with $c, D_{0}$ \emph{non-effective}
(i.e., guaranteed to exist by a compactness or ineffective-Siegel-style
argument but not numerically computable). Our pair is fully effective:
$c$ is computed from $c_{\mathrm{Brun}} / (2 \log 6)$ with explicit
$c_{\mathrm{Brun}} = 0.99$ (the Brun depth-$r = 10$ loss factor
\eqref{eq:delta-numeric}); $D_{0}$ is computed from the
binding BV constraint $D \geq e^{200}$. No invocation of Siegel zeros
or any ineffective input enters the derivation.

(ii) \textbf{Trade-off between $r$, $c$, and $D_{0}$.} The choice
$r = 11$ (Brun depth) is conservative. Smaller $r$ (e.g., $r = 6$)
gives larger combinatorial loss
$\delta_{r} \leq (e \cdot S_{1})^{r}/r!$ but smaller cumulative
remainder $R_{r}(z, Q) \leq Q^{r} \cdot \max|E_{V}|$, hence smaller
$D_{0}$. Specifically:
\begin{center}
\begin{tabular}{|c|c|c|c|}
\hline
$r$ & $\delta_{r}$ at $S_{1} = 1$ & approx.\ $D_{0}$
  (binding BV) & approx.\ $c$ \\
\hline
$6$ & $0.39$ & $e^{40} \approx 10^{17}$ & $0.085$ \\
$8$ & $0.10$ & $e^{100} \approx 10^{43}$ & $0.23$ \\
$10$ & $0.0060$ & $e^{150} \approx 10^{65}$ & $0.27$ \\
$11$ & $0.00151$ & $e^{200} \approx 10^{87}$ & $0.28$ \\
\hline
\end{tabular}
\end{center}
The choice $r = 11$ trades a larger $D_{0}$ for a sharper $c$. The
existential closure of EG\#203 (Lean axiom) only requires \emph{any}
effective pair; for the Lean axiom we can take $r = 6$ if desired,
giving $(c, D_{0}) \approx (0.085, 10^{17})$, which is much closer to
the empirical record.

(iii) \textbf{Pathways to further reduction.} A sharper analysis
using Bourgain--Glibichuk--Konyagin~\cite{BourgainGlibichukKonyagin2006}
\textbf{[EXISTS-LOCATION]} subgroup exp-sum bounds would give
power-saving in the per-prime discrepancy (replacing the
Erd\H{o}s--Tur\'an--Koksma logarithmic factors by power-of-$h_{q}$
saving), which would push the BV-binding constraint down to
$D_{0} \approx e^{50} \approx 10^{22}$. This is still well above
empirical but the structure of the proof would not change.


%==================== \S8 LEAN INTEGRATION ====================

\section{Lean integration}\label{sec:lean}

This preprint backs a single Lean 4 axiom in the kernel-verified
closure of Erd\H{o}s-Graham~\#203 at \cite{WilderEG203Lean2026}. In
plain math, the axiom asserts:
\[
  \exists\, c > 0,\ \exists\, D_{0} \in \mathbb{N} : \text{ for every
  ordinary } m, \text{ for every } D \geq D_{0},
  \quad c \cdot \mathfrak{S}(m) \cdot D \leq \pi_{V}(m, D).
\]
In Lean 4 source (the literal axiom statement uses Unicode
notation; the math is the displayed inequality above):

\begin{verbatim}
axiom wilder_2026_V_family_rosser_iwaniec :
  Exists (fun (c : Real) => Exists (fun (D0 : Nat) =>
    0 < c /\
      forall (m : Nat), Ordinary m -> forall (D : Nat), D0 <= D ->
        c * bateman_horn_singular_series_V m * (D : Real)
          <= (primeCountInBox m D : Real)))
\end{verbatim}

(Lean 4 source uses Unicode \texttt{\textbackslash exists}, \texttt{R},
\texttt{N}, \texttt{D\_0}, \texttt{forall}, \texttt{le} infix glyphs;
shown here in ASCII to avoid encoding artifacts.) The names:
\begin{itemize}
  \item \texttt{Ordinary m} is the predicate $\gcd(m, 6) = 1$.
  \item \texttt{bateman\_horn\_singular\_series\_V m} is the
    real-valued $\mathfrak{S}(m)$ from \cref{eq:S-defn-repeat}.
  \item \texttt{primeCountInBox m D} is the natural-number-valued
    $\pi_{V}(m, D)$ from \cref{eq:pi-V}.
\end{itemize}

\paragraph{Witness for the axiom.}
The pair $(c, D_{0}) = (0.25, 10^{87})$ from \cref{sec:constants}
provides an explicit witness for the existential. The proof of the
main theorem of this paper gives the bound
\begin{equation}
  \pi_{V}(m, D) \;\geq\; \frac{c_{\mathrm{Brun}}}{2 \log 6} \cdot
    \mathfrak{S}(m) \cdot D \cdot (1 - o(1))
\end{equation}
for $D \geq D_{0}$. Taking $c = 0.25 \leq c_{\mathrm{Brun}} /
(2 \log 6) \cdot 0.9$ and $D_{0} = 10^{87}$ satisfies the axiom
hypothesis.

\paragraph{Composition into the Erd\H{o}s-Graham~\#203 closure.}
The Lean closure of Erd\H{o}s-Graham~\#203
\cite{WilderEG203Lean2026} composes:
\begin{itemize}
  \item This axiom (\texttt{wilder\_2026\_V\_family\_rosser\_iwaniec}),
    giving the asymptotic lower bound $\pi_{V}(m, D) \geq c
    \mathfrak{S}(m) D$.
  \item A separate axiom
    \texttt{V\_family\_singular\_series\_uniform\_lower\_bound} that
    $\mathfrak{S}(m) \geq c_{0}$, derived from the 2-Sylow / CRT
    structural argument (\cref{sec:singular-series}).
  \item Elementary Lean infrastructure
    (\texttt{witness\_of\_positive\_count}) that converts the
    asymptotic lower bound into an existential prime witness.
\end{itemize}

The composition gives
\begin{verbatim}
theorem eg203_closed_atomic : EG203Closed
\end{verbatim}
which is kernel-verified. Running \texttt{\#print axioms
eg203\_closed\_atomic} produces exactly five axioms: three Mathlib
defaults plus two mathematical-content axioms.
\begin{verbatim}
[propext, Classical.choice, Quot.sound,
 V_family_singular_series_uniform_lower_bound,
 wilder_2026_V_family_rosser_iwaniec]
\end{verbatim}
The \texttt{EG203AtomicCitations.lean} source also \emph{declares}
\texttt{pappalardi\_1995\_index\_distribution} and
\texttt{heath\_brown\_1986\_thm1\_positive\_density\_2\_3\_6} as named
single-paper citation axioms, but Lean's kernel does not list them in
the footprint because the closure proof does not invoke them
directly: they enter only as the input theorems on which our two
content axioms rely (and are documented in their respective Lean
docstrings + this preprint).

The role of this preprint is to back the
\texttt{wilder\_2026\_V\_family\_rosser\_iwaniec} axiom with a
publishable mathematical argument. The companion axiom
\texttt{V\_family\_singular\_series\_uniform\_lower\_bound} follows
from the 2-Sylow / CRT structural argument
(\cref{sec:singular-series}) and is documented in detail in
\cite{WilderChain103}.

%==================== \S9 EMPIRICAL EVIDENCE ====================

\section{Empirical evidence}\label{sec:empirical}

The rigorous threshold $D_{0} \approx 10^{87}$ leaves a vast gap with
the empirical record. We document the latter for context.

\paragraph{Lean kernel-verified.}
Every ordinary $m \leq 200{,}000$ has a prime witness $(k, \ell)$
with $k + \ell \leq 12$. This is verified by the Lean file
\texttt{EG203BoundedClosure.lean} via \texttt{native\_decide} on a
search over $T_{12}$ for each $m$. Build SHA pinned at
\cite{WilderEG203Lean2026}.

\paragraph{Empirical sweeps (Python + sympy).}
We have performed empirical sweeps covering:
\begin{itemize}
  \item $m \in [1, 10^{6}]$: max witness diagonal $D = 17$.
  \item $m \in [1, 10^{7}]$: max witness diagonal $D = 19$.
  \item $m \in [1, 10^{8}]$: max witness diagonal $D = 19$.
  \item $m \in [1, 10^{9}]$: max witness diagonal $D = 22$.
  \item $m \in [10^{9}, 2 \cdot 10^{9}]$: max witness diagonal $D = 22$.
\end{itemize}
In all sweeps $\pi_{V}(m, D) \geq 1$ holds for all ordinary $m$ in the
range. SHA-pinned receipts at \cite{WilderEG203MasterLog}.

\paragraph{The empirical witness diagonal.}
The hardest $m$ at scale $10^{9}$ has witness diagonal $D = 22$. The
ratio $D / \log m$ stays around $1.0 \pm 0.2$ across all scales in the
sweep range, suggesting the true worst-case growth rate of
$D_{\min}(m) := \min\{D : \pi_{V}(m, D) \geq 1\}$ is $D_{\min}(m) \sim
\log m$ asymptotically. The rigorous bound proved here is
$D_{\min}(m) \leq D_{0}$ with $D_{0} \approx 10^{87}$ an absolute
constant in $m$; this is therefore much weaker than the empirical
$D_{\min}(m) \sim \log m$ functional dependence, but the rigorous
bound is sufficient for the existential closure of EG\#203 (any finite
$D_{0}$ suffices).

%==================== \S10 LIMITATIONS ====================

\section{Limitations and verification status}\label{sec:limitations}

\textbf{Expert review is solicited.} We list the specific points that
need verification by an expert in analytic number theory.

\paragraph{L1 --- Pappalardi-style quantitative input
(\cref{sec:kappa}).}
The proof of $\kappa_{V} = 0$ uses a quantitative rank-$2$
index-distribution bound: $\#\{p \leq z : H_{p} < (\log z)^{A}\} =
o(z / (\log z)^{B})$ for any $A, B > 0$. We attribute this to
Pappalardi 2003 via the reduction $H_{p} \geq \max(d_{p}(2),
d_{p}(3))$ + cyclic Pappalardi 1996. We have NOT verified the exact
exponent in Pappalardi 2003 verbatim; we have read the abstract and
theorem statement only.

\textbf{Action needed:} verify Pappalardi 2003 Theorem 1 (or the
rank-$2$ statement in \cite{PappalardiSusa2010}) gives the quantitative
form we claim. Alternatively, identify the strongest unconditional
rank-$2$ bound currently in the literature and adjust the argument
accordingly.

\paragraph{L2 --- Brun pure sieve at depth $r = 11$ (\cref{sec:sieve}).}
Proposition~\ref{prop:brun-pure} uses the Brun pure sieve at depth
$r = 11$, with explicit loss factor $\delta_{11} \leq 1.51 \times
10^{-3}$ derived from the Stirling form
$\delta_{r} \leq (e S_{1})^{r}/r!$ \eqref{eq:brun-stirling} and the
conservative bound $S_{1} \leq 1$ on the sum $\sum_{p} g_{V}(p)$. The
Stirling form is standard; the explicit reference is
Halberstam-Richert \S2.3 Lemma~2.2 + \S6 Thm.~6.1 (or
Iwaniec-Kowalski \S6.4 Prop.~6.7 + eq.~(6.7)). Verbatim line-checking
of these references has not been done.

\textbf{Action needed:} verify HR Lemma 2.2 + Thm 6.1 (or
IK Prop 6.7 + eq.~(6.7)) gives the precise Brun loss bound at depth
$r = 11$ with the constants we use.

\paragraph{L3 --- Cumulative remainder via per-prime ETK
(\cref{sec:bv}).}
The cumulative remainder $\sum_{d \leq Q} |E_{V}(D, d, 0)|$ at
squarefree level is bounded in Proposition~\ref{prop:BV} directly via
the per-prime Erd\H{o}s--Tur\'an--Koksma discrepancy
\eqref{eq:per-q-sharp} plus $\kappa_{V} = 0$ summation; this does
\emph{not} invoke the Iwaniec--Kowalski large-sieve cumulative
remainder (IK Thm 11.7). The 2D ETK inequality
(Drmota--Tichy~\cite[Thm.~1.21]{DrmotaTichy1997}, with absolute
constant $C_{2} = 4$) and the standard 1D logarithmic-sum bound
(\cite[Lemma~3.1]{DrmotaTichy1997}) are the load-bearing inputs;
neither is exotic, both are textbook-standard.

\textbf{Action needed:} verify Drmota-Tichy 1997 Theorem 1.21
(absolute constant $C_{s} = 4$ in dim.\ $s = 2$) and Lemma 3.1
verbatim.

\paragraph{L4 --- Asymptotic sieve at $\kappa = 0$
(\cref{sec:almost-prime}).}
The asymptotic prime count $\pi_{V}(m, D) \sim |T_{D}| \cdot
\mathfrak{S}(m) / \log(m \cdot 6^{D})$ at $\kappa_{V} = 0$ is invoked
in \S\ref{sec:almost-prime} for the conversion from Brun-pure-sieve
lower bound to prime count. The asymptotic-sieve framework at
$\kappa = 0$ is covered in
Halberstam--Richert \cite[\S25]{HalberstamRichert1974} and
Iwaniec--Kowalski \cite[\S22]{IwaniecKowalski2004}. We did not verify
the verbatim statement of the $\kappa = 0$ asymptotic in either
reference. As noted in Remark~\ref{rem:almost-prime-final}, an
alternative explicit derivation via the iterated Brun upper sieve
gives the same conclusion with slightly worse constants.

\textbf{Action needed:} verify the $\kappa = 0$ asymptotic in
HR \S25 or IK \S22; alternatively, write out the iterated Brun upper
sieve derivation in full.

\paragraph{L5 --- Sylow / CRT density argument
(\cref{sec:singular-series}).}
The lower bound $\mathfrak{S}(m) \geq 1/256$ depends on the
structural 2-Sylow / orbit-stabilizer / CRT argument of
\cite{WilderChain103}. This argument is original to this project,
Lean-verified (via \texttt{Chain103UniversalDensity}), and
Python-cross-checked across three independent period extensions of
the base $(2520, 5040)$. The translation from ``uncovered base
density $\geq 1/128$'' to ``singular series $\geq 1/128 \cdot$
Mertens tail'' is standard but not formally Lean-checked.

\textbf{Action needed:} formal Lean proof of the singular-series
translation, or independent paper verifying the 2-Sylow uncovered
density argument.

\paragraph{L6 --- Effective constant $D_{0} = 10^{87}$.}
The threshold $D_{0}$ is derived from four constraints
(\S\ref{sec:constants}): (i) BV convergence, (ii) Brun-sieve threshold
$z = D^{1/3}$, (iii) the $\mathfrak{S}(m) \geq c_{0}$ lower bound
from \cite{WilderChain103}, (iv) the $W(z) \to \mathfrak{S}(m)$ rate
and asymptotic-sieve convergence. The \emph{binding} constraint is
(i) BV at $D_{1} \geq e^{200} \approx 10^{87}$. Constraints (ii),
(iii), (iv) are all non-binding by a wide margin.

A tighter analysis via Bourgain--Glibichuk--Konyagin 2006 subgroup
exp-sum power-saving could reduce $D_{0}$ to $\approx 10^{22}$. The
trade-off with the Brun depth $r$ is tabulated in
\S\ref{sec:constants}: smaller $r$ gives larger $\delta_{r}$ but
smaller $D_{0}$, e.g.\ $r = 6$ gives $(c, D_{0}) \approx (0.085,
10^{17})$.

\textbf{Action needed:} (optional) tighten via BGK 2006 to reduce
$D_{0}$, or use smaller $r$ for the existential closure.

\paragraph{L7 --- Verbatim citation verification.}
Each bibliography entry is tagged with one of three labels:
\begin{itemize}
  \item \texttt{[VERIFIED]}: paper read and theorem statement (or
    abstract) matches the citation. (7 entries: Bombieri 1965,
    Friedlander-Iwaniec 1998, Maynard 2019, Mauduit-Rivat 2010,
    Erd\H{o}s-Graham 1980, Erd\H{o}s-Pomerance 1985, Rosser-Schoenfeld
    1962.) Two of these (Iwaniec 1980, Heath-Brown 1986) are
    \texttt{[VERIFIED at paper existence + statement of theorem in
    use]}: paper exists, the cited theorem statement was read, but the
    verbatim wording of supporting lemmas was not line-checked.
  \item \texttt{[EXISTS-LOCATION]}: paper exists, theorem location is
    known but the verbatim line was not line-checked. (Used for the
    remaining literature citations: Halberstam-Richert 1974,
    Iwaniec-Kowalski 2004, Pappalardi 1996, Pappalardi 2003,
    Pappalardi-Susa 2010, Drmota-Tichy 1997, Greaves 2001, Hooley 1971,
    Brun 1915, Bourgain-Glibichuk-Konyagin 2006, Heath-Brown-Konyagin
    2000, Heath-Brown 2001, Vinogradov 1965.)
  \item \texttt{[NEW]}: this preprint / companion notes. Three entries
    (Wilder 2026a, 2026b, 2026c).
\end{itemize}

\textbf{Action needed:} systematic verbatim verification of all
[EXISTS-LOCATION] citations. The most load-bearing are
Halberstam-Richert \S2.3 + \S6 (L2), Drmota-Tichy Thm.~1.21 (L3),
Pappalardi 2003 / Pappalardi-Susa 2010 (L1), and the
Halberstam-Richert \S25 or Iwaniec-Kowalski \S22 asymptotic-sieve
chapter (L4).

\paragraph{L8 --- The empirical-rigorous gap.}
The gap from rigorous $D_{0} \approx 10^{87}$ to empirical $D \leq 22$
is not addressed. We conjecture but do not prove that the rigorous
$D_{0}$ can be reduced to $D_{0} = O(\log m)$ for some absolute
constant.

\textbf{Action needed:} (optional, for a stronger result) close the
gap by exploiting the specific structure of $\langle 2, 3 \rangle$
beyond what we use here.

\paragraph{L9 --- Preprint status.}
This is a preprint, not a refereed publication. The author is
responsible for the content; literature citations are tagged in L7
above with their verification status. Source files for the working
drafts of each section are preserved alongside this preprint in the
source directory.

%==================== BIBLIOGRAPHY ====================

\begin{thebibliography}{99}

\bibitem[Bombieri 1965]{Bombieri1965}
E.~Bombieri.
\emph{On the large sieve.}
Mathematika 12 (1965), 201--225.
\texttt{[VERIFIED at paper existence and abstract]}

\bibitem[Bourgain-Glibichuk-Konyagin 2006]{BourgainGlibichukKonyagin2006}
J.~Bourgain, A.~Glibichuk, S.~V.~Konyagin.
\emph{Estimates for the number of sums and products and for
exponential sums in fields of prime order.}
J. London Math. Soc. (2) 73 (2006), no.~2, 380--398.
\texttt{[EXISTS-LOCATION]}

\bibitem[Brun 1915]{Brun1915}
V.~Brun.
\emph{\"Uber das Goldbachsche Gesetz und die Anzahl der Primzahlpaare.}
Arch. f. Math. og Naturvid. B34 (1915), no.~8, 1--19.
\texttt{[EXISTS-LOCATION; classical reference]}

\bibitem[Drmota-Tichy 1997]{DrmotaTichy1997}
M.~Drmota, R.~F.~Tichy.
\emph{Sequences, Discrepancies and Applications.}
Lecture Notes in Math. 1651, Springer, 1997.
\texttt{[EXISTS-LOCATION; standard reference for ET inequality]}

\bibitem[Erd\H{o}s-Graham 1980]{ErdosGraham1980}
P.~Erd\H{o}s, R.~L.~Graham.
\emph{Old and new problems and results in combinatorial number theory.}
L'Enseignement Math\'ematique Monographies 28, Geneva, 1980.
Problem~203, p.~27.
\texttt{[VERIFIED]}

\bibitem[Erd\H{o}s-Pomerance 1985]{ErdosPomerance1985}
P.~Erd\H{o}s, C.~Pomerance.
\emph{On the normal number of prime factors of $\phi(n)$.}
Rocky Mountain J. Math. 15 (1985), 343--352.
\texttt{[VERIFIED at paper existence; the bound we use is from the
Pomerance 1985 companion paper "On the distribution of multiplicative
orders"]}

\bibitem[Friedlander-Iwaniec 1998]{FriedlanderIwaniec1998}
J.~Friedlander, H.~Iwaniec.
\emph{The polynomial $X^{2} + Y^{4}$ captures its primes.}
Ann. of Math. (2) 148 (1998), no.~3, 945--1040.
\texttt{[VERIFIED at title and abstract; key prior art for
asymptotic-formula sieves at $\kappa = 1/2$]}

\bibitem[Greaves 2001]{Greaves2001}
G.~Greaves.
\emph{Sieves in number theory.}
Ergebnisse der Mathematik 43, Springer, 2001.
\texttt{[EXISTS-LOCATION; standard reference for sieve functions at
various $\kappa$]}

\bibitem[Halberstam-Richert 1974]{HalberstamRichert1974}
H.~Halberstam, H.-E.~Richert.
\emph{Sieve methods.}
London Math. Soc. Monographs 4, Academic Press, 1974.
\texttt{[EXISTS-LOCATION; standard reference]}

\bibitem[Heath-Brown 1986]{HeathBrown1986}
D.~R.~Heath-Brown.
\emph{Artin's conjecture for primitive roots.}
Quart. J. Math. Oxford Ser. (2) 37 (1986), no.~145, 27--38.
\texttt{[EXISTS-LOCATION; verified at paper existence + statement of
positive-density result]}

\bibitem[Heath-Brown 2001]{HeathBrown2001}
D.~R.~Heath-Brown.
\emph{Primes represented by $x^{3} + 2 y^{3}$.}
Acta Math. 186 (2001), no.~1, 1--84.
\texttt{[EXISTS-LOCATION]}

\bibitem[Heath-Brown-Konyagin 2000]{HeathBrownKonyagin2000}
D.~R.~Heath-Brown, S.~V.~Konyagin.
\emph{New bounds for Gauss sums derived from kth powers, and for
Heilbronn's exponential sum.}
Quart. J. Math. Oxford 51 (2000), 221--235.
\texttt{[EXISTS-LOCATION]}

\bibitem[Hooley 1971]{Hooley1971}
C.~Hooley.
\emph{On the Brun-Titchmarsh theorem.}
J. Reine Angew. Math. 255 (1971), 60--79.
\texttt{[EXISTS-LOCATION; for Brun-Hooley sieve improvement]}

\bibitem[Iwaniec 1980]{Iwaniec1980}
H.~Iwaniec.
\emph{Rosser's sieve.}
Acta Arith. 36 (1980), 171--202.
\texttt{[VERIFIED at paper existence + general framework; the $\kappa
= 0$ specialization we cite is not the explicit subject of this
paper]}

\bibitem[Iwaniec-Kowalski 2004]{IwaniecKowalski2004}
H.~Iwaniec, E.~Kowalski.
\emph{Analytic Number Theory.}
AMS Colloquium Publications 53, 2004.
\texttt{[EXISTS-LOCATION; Theorems 6.1 (Brun-Hooley), 11.6
(Rosser-Iwaniec linear sieve), 11.7 (large sieve), 22.x (asymptotic
sieve)]}

\bibitem[Maynard 2019]{Maynard2019}
J.~Maynard.
\emph{Primes with restricted digits.}
Invent. Math. 217 (2019), no.~1, 127--218.
\texttt{[VERIFIED at title and abstract]}

\bibitem[Mauduit-Rivat 2010]{MauduitRivat2010}
C.~Mauduit, J.~Rivat.
\emph{Sur un probl\`eme de Gelfond: la somme des chiffres des nombres
premiers.}
Ann. of Math. (2) 171 (2010), no.~3, 1591--1646.
\texttt{[VERIFIED at title and abstract]}

\bibitem[Pappalardi 1996]{Pappalardi1995}
F.~Pappalardi.
\emph{On the order of finitely generated subgroups of
$\mathbb{Q}^{\times}$ (mod $p$) and divisors of $p - 1$.}
J. Number Theory 57 (1996), no.~2, 207--222.
(Submitted 1995, published 1996; bib key
\texttt{Pappalardi1995} retained for backward compatibility.)
\texttt{[EXISTS-LOCATION; Theorem 1 statement read at abstract level]}

\bibitem[Pappalardi 2003]{Pappalardi2003}
F.~Pappalardi.
\emph{Square-free values of the order function.}
J. Number Theory 100 (2003), 251--263.
\texttt{[EXISTS-LOCATION; rank-2 quantitative refinement]}

\bibitem[Pappalardi-Susa 2010]{PappalardiSusa2010}
F.~Pappalardi, A.~Susa.
\emph{An analogue of Artin's conjecture for multiplicative subgroups
of the rationals.}
Arch. Math. 95 (2010), no.~1, 27--37.
\texttt{[EXISTS-LOCATION]}

\bibitem[Pomerance 1985]{Pomerance1985}
C.~Pomerance.
\emph{On the distribution of multiplicative orders.}
1985 (informal note / appendix to Erd\H{o}s-Pomerance 1985).
\texttt{[EXISTS-LOCATION; the bound $\#\{p \leq z : \mathrm{ord}_{p}(a)
< z^{1-\alpha}\} \ll z^{1 - \alpha/2}$ is the standard
Erd\H{o}s-Pomerance bound]}

\bibitem[Rosser-Schoenfeld 1962]{RosserSchoenfeld1962}
J.~B.~Rosser, L.~Schoenfeld.
\emph{Approximate formulas for some functions of prime numbers.}
Illinois J. Math. 6 (1962), 64--94.
\texttt{[VERIFIED at paper and Theorem 23 statement; used for Mertens
explicit constants]}

\bibitem[Wilder 2026a]{WilderChain103}
J.~Wilder.
\emph{Exact structural characterization of uncovered cells in the
$(2520, 5040)$ period base for the V-family.}
Companion note to the present preprint, 2026. Includes a Lean 4
formalization of the universal density statement
(\texttt{Chain103UniversalDensity.lean}) and Python verification
across three independent period extensions.
\texttt{[NEW; Lean-formalized; awaiting peer review]}

\bibitem[Wilder 2026b]{WilderEG203Lean2026}
J.~Wilder.
\emph{Erd\H{o}s-Graham~\#203: a kernel-verified closure in Lean 4.}
Companion Lean 4 development, 2026.
\texttt{[NEW; Lean-formalized; awaiting peer review]}

\bibitem[Wilder 2026c]{WilderEG203MasterLog}
J.~Wilder.
\emph{Erd\H{o}s-Graham~\#203: a verification log.}
Companion technical note recording the empirical sweeps and
verification chains referenced in \S\ref{sec:empirical}, 2026.
\texttt{[NEW]}

\bibitem[Vinogradov 1965]{Vinogradov1965}
A.~I.~Vinogradov.
\emph{The density hypothesis for Dirichlet $L$-series.}
Izv. Akad. Nauk SSSR Ser. Mat. 29 (1965), 903--934.
\texttt{[EXISTS-LOCATION; classical BV co-discoverer reference]}

\end{thebibliography}

\end{document}
